% Modernised reading — fonds Alexandre Grothendieck, shelfmark 159,
% pages 1-33 (the whole folder).
%
% This is an INTERPRETATION, not a transcription. It re-reads the pages in
% today's notation and vocabulary, states the hypotheses the manuscript leaves
% implicit, and drops the critical apparatus so that the mathematics reads
% straight through. Everything the brackets and underlines carried — a doubtful
% reading, a notation he used differently, a missing passage — moves into a
% footnote.
%
% It is held to being mathematically correct as it stands, not merely faithful
% to the page. Where the manuscript is loose or mistaken, the true statement is
% what appears here, and the footnote says what the page has. In any dispute of
% substance the transcription governs: whoever cites Grothendieck must cite
% batch-01.fr.tex and batch-02.fr.tex.
%
% Pass: Opus 5 (claude-opus-5), 2026-09-19 — first pass, unchecked against the
% pages by a human. The model was named explicitly in the request; it is the
% model this reading actually ran on.
%
% Scope: a SINGLE document for the whole folder. The shelfmark holds two
% batches — pages 1-20 and pages 21-33 — read together in one pass, and no part
% of the folder is read in another file. Page 13 carries nothing but the
% archivists' number and page 23 is a blank leaf; both are absent from the
% transcriptions and from here.
%
% Notation fixed for the whole document, and departing from the pages:
%   - H = PGL_2 = Aut(P^1), where the pages write PGL(1). His index counts the
%     RELATIVE DIMENSION of the Brauer-Severi scheme, not the rank of the
%     module: page 1 writes PGL(r)_S for a P^r-bundle. The convention is
%     confirmed internally by page 15, where chi_DR(X,g) = 3 chi_DR(X,O_X)
%     requires rk(g) = 3 = dim PGL_2.
%   - O_S, O_X for his underlined O.
%   - z (fraktur z) for the raised glyph the transcriber of pages 1-20 could
%     not name. It is an EDITORIAL PLACEHOLDER, not a reading: see the note in
%     the spine section and the footnote at page 14.
%   - omega = Omega^1_{X/S} on a relative curve; Omega for the invertible (or
%     locally free) module with J/J^2 = Omega, identified with Omega^1_{X/S}
%     when a projective connection is in play; Omega-check for its dual.
%   - varpi_i for the coefficients of C restricted to Omega, as the
%     transcription of pages 27-33 sets them.
%   - M_g^z for the moduli of pairs (curve, projective structure) over M_g;
%     M_g^z(X) for the character variety of pi = pi_1(X,x) on pages 16-17.
%     The two are DIFFERENT spaces of the same dimension 6g-6, and the pages
%     use one symbol for both. See the conventions.
%
% Substantive departures from the pages, each footnoted where it occurs:
%   - page 8: P^infty_{X/S} is an O_X-algebra, so the power series ring is
%     O_X[[delta f_1,...,delta f_n]], not O_S[[...]] as the page writes.
%   - page 14: a) is strictly stronger than b) as the three conditions stand;
%     what is proved there is b) <=> c) and a) => b). The equivalence a) <=> b)
%     for a curve of genus >= 2 carrying a projective connection follows from
%     the theorem of pages 21-22, which the folder does not invoke here.
%   - page 15: the relative dimension of M_g^z over M_g is 3g-3, not 6g-6.
%     The leaf is not repaired; see the section devoted to it.
%   - page 21: the two cases treated (Borel, normaliser of a maximal torus) do
%     not exhaust the proper algebraic subgroups of PGL_2 in characteristic 0 —
%     A_4, S_4, A_5 remain. The same device the page uses settles them.
%   - page 30: the first term of C(omega) is written varpi_{-1}, which page 27
%     excludes; the term is not read with certainty and is not reconstructed.
%
% What the folder announces and never establishes: the CANONICAL projective
% connection its own title promises (page 8 asks the question and leaves it);
% the characterisation of the Teichmueller locus inside the character variety
% (page 17 states it as a supposition); the injectivity of M_ann(X) ->
% M_g^z(X), asked on page 17 and not answered; item b) of the theorem of
% page 25, which breaks off mid-sentence; item c) of page 26, almost entirely
% lost; the whole of page 9 after the word "Posant".
\documentclass[11pt,a4paper]{article}
% Le préambule commun à toutes les transcriptions du fonds Grothendieck.
%
% Il tient en une idée : **l'incertitude se compose, elle ne se résout pas.**
% Une transcription qui lisse un mot douteux a détruit la seule chose pour
% laquelle on la faisait — un lecteur doit pouvoir savoir, à chaque endroit,
% ce qui était sur le feuillet et ce qui a été ajouté. Les six macros qui
% suivent sont donc l'appareil critique, et non de la décoration.
%
% Les mêmes commandes sont reconnues par `scripts/render.mjs`, qui produit la
% vue de lecture du site. Une macro ajoutée ici doit l'être aussi là-bas,
% faute de quoi le rendu échouera bruyamment — ce qui est voulu : un rendu
% silencieusement faux est pire qu'un rendu absent.

\ProvidesPackage{grothendieck}[2026/08/08 Transcriptions du fonds Grothendieck]

\usepackage{amsmath,amssymb,amsthm}
\usepackage{mathrsfs}
\usepackage{stmaryrd}
% Les diagrammes commutatifs sont partout dans ce fonds : c'est la langue
% même de la géométrie algébrique de Grothendieck, et une transcription qui
% les rendrait en prose ne transcrirait rien.
\usepackage{tikz-cd}
% Français par défaut — c'est la langue des feuillets — mais l'anglais est
% chargé aussi : la traduction et le résumé sont des documents anglais, et la
% typographie française (espaces avant les deux-points) y serait une faute.
% Ces documents basculent par \selectlanguage{english} en tête de corps.
\usepackage[english,french]{babel}
\usepackage{fontspec}
\usepackage{geometry}
\geometry{a4paper,margin=25mm,marginparwidth=28mm}
\usepackage{marginnote}
\usepackage{xcolor}
% ulem, not soul. `soul`'s \st and \ul are built on a character-by-character
% scan that XeLaTeX's font machinery breaks: the document fails with an
% undefined control sequence rather than degrading. ulem does the same two jobs
% and is Unicode-engine safe. `normalem` stops it hijacking \emph, which the
% transcriptions use for Grothendieck's own emphasis.
\usepackage[normalem]{ulem}
\usepackage{hyperref}
\hypersetup{colorlinks=true,linkcolor=black,urlcolor=black,pdfborder={0 0 0}}

% --- Métadonnées du lot -----------------------------------------------------
% Reprises telles quelles par le script de rendu, qui les affiche en tête de
% la vue de lecture. Elles fixent ce que le fichier prétend couvrir : sans
% elles, une transcription flotte sans point d'attache dans un fonds de seize
% mille feuillets.

\newcommand{\folder}[1]{\gdef\grothFolder{#1}}
\newcommand{\batch}[1]{\gdef\grothBatch{#1}}
\newcommand{\leaves}[2]{\gdef\grothFirst{#1}\gdef\grothLast{#2}}
\newcommand{\foldertitle}[1]{\gdef\grothFoldertitle{#1}}
\gdef\grothFolder{?}\gdef\grothBatch{?}\gdef\grothFirst{?}\gdef\grothLast{?}\gdef\grothFoldertitle{}

% --- Le résumé --------------------------------------------------------------
%
% Il ouvre la lecture modernisée, sous le titre « Résumé », et il est composé
% en petit. Ce n'est pas un ornement : c'est le seul passage du document qui
% ne soit pas des mathématiques, et un lecteur doit pouvoir le reconnaître
% avant de l'avoir lu — puis le sauter s'il connaît déjà le sujet, ou s'y
% arrêter s'il ne le connaît pas. Le filet qui le suit marque la reprise du
% corps.

\newenvironment{resume}
  {\small\setlength{\parskip}{0.45em}}
  {\par\medskip\hrule\medskip}

% --- Mots-clés ---------------------------------------------------------------
%
% En anglais, en clôture du résumé de la lecture modernisée : le vocabulaire
% moderne sous lequel chercher ce que les feuillets construisent. Le manifeste
% du site (`npm run manifest`) les extrait pour étiqueter la cote entière —
% cette ligne est la source unique des tags, il n'y a pas de fichier à côté.
\newcommand{\keywords}[1]{\par\smallskip\noindent
  {\footnotesize\textsc{Keywords} --- \textit{#1}}\par}

% --- L'appareil critique ----------------------------------------------------

% Le numéro de feuillet, en marge. C'est l'ancre qui permet de régler un
% désaccord sur une formule en regardant *un* feuillet plutôt qu'une chemise
% de six cents. Le site s'en sert aussi pour tourner les pages du fac-similé
% quand on fait défiler la transcription.
\newcommand{\leaf}[1]{%
  \marginnote{\footnotesize\color{gray!70}\textsf{#1}}%
  \label{leaf:#1}\ignorespaces}

% Illisible : on ne devine pas. Un mot inventé qui se lit comme les autres est
% le pire résultat possible.
\newcommand{\ill}{\textcolor{red!60!black}{[\,\ldots\,]}}

% Lecture douteuse : le mot est donné, et signalé comme incertain.
\newcommand{\uncertain}[1]{\uline{#1}}

% Ajout de l'éditeur — une lettre manquante, un mot sous-entendu.
\newcommand{\add}[1]{\textcolor{blue!50!black}{[#1]}}

% Biffé par Grothendieck lui-même. Ce qu'il a rayé fait partie du document :
% une rature dit souvent où la pensée a changé de direction.
\newcommand{\struck}[1]{\sout{#1}}

% Note du transcripteur, sur l'état matériel du feuillet.
\newcommand{\note}[1]{\par\smallskip{\small\color{gray!60!black}\textit{[#1]}}\par\smallskip}

% Note marginale de Grothendieck, distincte du corps du texte.
\newcommand{\marginal}[1]{\par\smallskip\noindent\hspace{1em}%
  {\small\color{gray!60!black}#1}\par\smallskip}

% --- Titre ------------------------------------------------------------------

\AtBeginDocument{%
  \begin{center}
    {\large\bfseries Fonds Alexandre Grothendieck --- cote n\textsuperscript{o}~\grothFolder}\\[2pt]
    {\itshape\grothFoldertitle}\\[4pt]
    {\small feuillets \grothFirst--\grothLast\ (lot \grothBatch)}\\[2pt]
    {\footnotesize\color{gray!60!black}Transcription établie d'après le fac-similé numérisé
     par l'Université de Montpellier.\\
     Les lectures incertaines sont soulignées, les illisibles notées [\,\ldots\,],
     les ajouts entre crochets.}
  \end{center}
  \vspace{1em}}

\endinput


\folder{159}
\pages{1}{33}
\dating{s.d. — le groupe « Dérivateurs » (157-1 à 160) est daté 1990-[1991]}
\watermark{Édition de démonstration\\interprétation personnelle de l'œuvre}
\foldertitle{Connexions projectives et uniformisation : notes manuscrites (s.d.) — lecture modernisée du dossier entier}

\begin{document}

\section*{Résumé}

\begin{resume}
Sur une surface, on sait ce qu'est une carte : un morceau du plan, collé à
d'autres morceaux du plan. Ce qui fait la géométrie, c'est la règle de
recollement — ce qu'on autorise comme changement de carte. Si l'on n'autorise
que les déplacements, on fait de la géométrie euclidienne ; si l'on autorise
toutes les bijections différentiables, on ne fait plus de géométrie du tout,
seulement de la topologie. Entre les deux, il y a des choix intéressants. Une
\emph{structure projective} sur une surface, c'est le choix suivant : les
cartes sont des morceaux de la droite projective complexe $\mathbf{P}^1$ — la
sphère de Riemann — et les changements de carte sont les seules
transformations $z \mapsto (az+b)/(cz+d)$, celles qui préservent la structure
projective de la sphère.

Ce choix a une vertu particulière, et c'est elle qui fait tout le sujet. Les
transformations $z \mapsto (az+b)/(cz+d)$ forment un groupe, $\mathrm{PGL}_2$,
et une carte projective est rigide : deux cartes qui coïncident sur un ouvert
coïncident partout où elles sont définies toutes deux. On peut donc, en
partant d'une carte et en la prolongeant le long des chemins, fabriquer une
seule application de la surface \emph{universelle} — le revêtement universel
$\widetilde{X}$ — vers $\mathbf{P}^1$. On l'appelle l'\emph{application
développante}. Elle n'est pas définie sur $X$ mais sur $\widetilde{X}$, et le
défaut se mesure par un homomorphisme $\rho : \pi_1(X) \to \mathrm{PGL}_2$,
la \emph{monodromie} : faire le tour d'un lacet change la développante par
la transformation $\rho(\gamma)$.

L'uniformisation, c'est le cas privilégié où tout cela se referme. Si
l'application développante est un isomorphisme de $\widetilde{X}$ sur un
disque $D \subset \mathbf{P}^1$ et si la monodromie identifie $\pi_1(X)$ à un
groupe de transformations agissant librement sur $D$ avec quotient compact,
alors $X$ \emph{est} $D/\pi_1(X)$ : la surface a été construite à partir d'un
morceau du plan et d'un groupe. C'est le théorème d'uniformisation, et c'est
une structure projective très particulière. Toutes les autres structures
projectives sur la même surface sont des voisines de celle-là, et la question
naturelle est : combien y en a-t-il, comment sont-elles rangées, et
reconnaît-on celle qui uniformise ?

Le dossier attaque cette question par un autre bout que l'analyse complexe,
et c'est ce qui le rend intéressant. Il veut une théorie algébrique, valable
au-dessus d'une base quelconque et sans convergence ni revêtement universel :
là où l'analyste a des cartes, l'algébriste n'a que des voisinages
infinitésimaux. Une structure projective devient alors une
\emph{connexion projective}, qui est la donnée de deux choses : un fibré en
droites projectives $P$ sur la courbe $X$, avec une section — ce qui, par un
dictionnaire que les six premières pages établissent soigneusement, revient à
une extension de modules $0 \to \Omega \to E \to \mathcal{O}_X \to 0$ — et une
identification, ordre par ordre, entre le voisinage formel de la section dans
$P$ et le voisinage formel de la diagonale dans $X \times X$. La seconde donnée
est la traduction algébrique exacte de « une carte à valeurs dans
$\mathbf{P}^1$ » : elle dit que $X$ ressemble à $\mathbf{P}^1$ à tout ordre
fini, et de façon cohérente.

De là le dossier tire une série de comptes. Sur une courbe de genre $g$, les
connexions projectives forment un espace affine dont les translations sont les
\emph{différentielles quadratiques}, les sections de $\omega^{\otimes 2}$ ;
il y en a $3g-3$ en dimension, exactement le nombre de paramètres dont dépend
la courbe elle-même. L'espace des couples (courbe, structure projective) a
donc $6g-6$ paramètres, et c'est aussi la dimension de l'espace des
monodromies possibles à conjugaison près — une coïncidence qui n'en est pas
une, et qui est le contenu géométrique de tout l'exercice. Un autre compte,
plus classique, apparaît au passage : la différence entre les
$g(g+1)/2$ produits de deux formes différentielles et les $3g-3$
différentielles quadratiques vaut $\tfrac12 (g-2)(g-3)$, qui est la dimension de
l'espace des quadriques passant par la courbe canonique.

Le résultat le plus net du dossier est ailleurs, et il est de nature
négative. Sur une courbe de genre au moins $2$, aucune connexion projective ne
peut être « réductible » : le groupe $\mathrm{PGL}_2$ ne se ramène jamais à un
sous-groupe plus petit. La démonstration tient en six lignes de calcul de
degrés — un sous-fibré invariant aurait le degré $g-1$, ce qui force une
section nulle et une contradiction — et sur le corps des complexes elle dit
que l'image de la monodromie est Zariski-dense. Autrement dit : une structure
projective algébrique est toujours, en un sens précis, aussi peu dégénérée que
possible. C'est le genre d'énoncé qui rend un espace de modules utilisable, et
c'est aussi, dans le dossier, la pièce qui manquait ailleurs.

Le dernier tiers des feuillets — dix pages, reprises cinq ou six fois — est
d'une autre facture. Il ne parle plus de courbes du tout, mais d'algèbres :
une algèbre complète, augmentée, dont le premier gradué est inversible, et
munie d'une dérivation compatible. La question est de savoir si une telle
algèbre est toujours isomorphe au modèle, l'algèbre des séries formelles, et
de combien de façons. La réponse est aussi rigide qu'on peut l'espérer :
l'isomorphisme existe, il est unique dès qu'on l'a fixé à l'ordre $3$, et
l'ambiguïté à cet ordre est exactement une différentielle quadratique. C'est
la démonstration technique de tout ce qui précède, écrite après coup et pour
elle-même.

Un mot sur ce qui n'est pas là. Le titre de la première page est de sa main et
promet des connexions projectives \emph{canoniques} ; le dossier n'en construit
aucune, et la page 8 pose la question — une connexion projective est-elle
canoniquement déterminée par la courbe ? — sans y répondre. L'uniformisation
proprement dite n'est pas atteinte non plus : la page 17 décrit le lieu des
structures qui uniformisent, et introduit la description par le mot
« on suppose ». Ce que le dossier livre est l'appareil, non le théorème. Il
faut le dire, parce que le titre promet davantage.

Un mot enfin sur ce que ces feuillets ne disent pas. Le dossier est rangé dans
un groupe que les archivistes datent de 1990-[1991] et intitulent
« Dérivateurs » ; rien de ce qui est écrit sur ces trente-trois pages ne relève
des dérivateurs. C'est le vocabulaire des années EGA et SGA. C'est un constat
de contenu, refaisable sur la transcription, et il ne date rien : un
mathématicien revient sur un sujet ancien quand il lui plaît, et un dossier
d'archives est une chemise de papier. Les deux transcripteurs ajoutent des
observations matérielles — du papier à portées musicales imprimées aux pages
30 et 32-33, du listing informatique en accordéon ailleurs — mais ce sont des
récupérations qu'il a pratiquées sur des décennies, faites ici sur un
fac-similé, et elles ne tranchent rien non plus. La datation portée en tête
reste celle de Montpellier, inchangée.

Les noms modernes sous lesquels chercher : structure projective sur une
surface de Riemann, connexion projective et dérivée schwarzienne, opers au
sens de Beilinson et Drinfeld, application développante et application de
monodromie, variété des caractères, espace de Teichmüller, différentielles
quadratiques, fibrés de Brauer-Severi, parties principales et algèbres de
jets, connexion de Gauss-Manin et application de Kodaira-Spencer, théorème de
Max Noether sur les quadriques de la courbe canonique.

\keywords{projective structure, projective connection, oper, Schwarzian derivative, developing map, monodromy representation, uniformization, Fuchsian group, Teichmuller space, character variety, Zariski density, irreducible local system, Brauer-Severi scheme, principal parts, jet algebra, formal completion, affine group, torsor, quadratic differential, canonical curve, Max Noether theorem, Gauss-Manin connection, Kodaira-Spencer map, de Rham cohomology, Hodge filtration, symplectic self-duality, moduli of curves, Lie algebra cohomology, Euler-Poincare characteristic, adjoint bundle, Riemann-Hilbert correspondence, complete augmented algebra, rigidity}
\end{resume}

\pagerange{1}{33}
\section*{Le fil du dossier, ses reprises, et les conventions}

Le dossier n'est pas un texte suivi. Il a un titre, de sa main, porté en tête
de la page 1 : « Connexions projectives canoniques et uniformisation des
courbes algébriques ». Il a ensuite deux grands massifs qui ne se rejoignent
jamais explicitement, et un théorème isolé entre les deux.

\pagerange{1}{33}
\subsection*{Les stations}

\begin{itemize}
\item \textbf{Pages 1 à 6} — le dictionnaire. Fibré de Brauer-Severi avec
section, réduction du groupe structural au groupe affine, extension
$0 \to \Omega \to E \to \mathcal{O}_S \to 0$ : trois descriptions d'un seul
objet. Puis le complété formel le long de la section, et le calcul explicite
de l'action de $\Omega$ dessus. La page 2 est un départ abandonné, repris de
zéro à la page 3 ; on la garde pour ce qu'elle est.
\item \textbf{Pages 7 à 9} — la définition. Une connexion projective sur
$X/S$, l'identification forcée $\Omega \simeq \Omega^1_{X/S}$, la condition à
l'ordre $2$, l'intégrabilité, et l'écriture en coordonnées locales. La
page 9 s'arrête au premier tiers, sur le mot « Posant ».
\item \textbf{Pages 10 à 12} — la dimension relative $1$. La cohomologie de
de Rham d'une famille de courbes, son autodualité symplectique, les
scindages de la filtration de Hodge, Gauss-Manin et Kodaira-Spencer, et le
compte $\tfrac12(g-2)(g-3)$.
\item \textbf{Pages 14 à 18} — les modules. L'espace $M_g^{\mathfrak{z}}$, le
critère de non-dégénérescence sous trois formes, le calcul de dimension par
la caractéristique d'Euler-Poincaré de de Rham, puis le tableau analytique :
variété des caractères, action du groupe modulaire, plongement de l'espace de
Teichmüller. La page 13 ne porte que le numéro des archivistes.
\item \textbf{Pages 19 et 20} — le calcul en coordonnées, qui aboutit à la
formule encadrée $u_\omega(c) = c + d^1_{\omega/S}(\omega)$.
\item \textbf{Pages 21 et 22} — le théorème de non-réduction, et la densité
Zariski de la monodromie sur $\mathbf{C}$. Ces deux pages achèvent un
argument commencé plus haut dans le dossier.
\item \textbf{Pages 24 à 33} — la théorie formelle, reprise cinq ou six fois.
C'est un second sujet, et c'est aussi la démonstration de ce que les pages 7
à 9 supposaient. La page 23 est un feuillet blanc.
\end{itemize}

\pagerange{24}{33}
\subsection*{Les reprises, et leur ordre}

Le dernier massif n'est pas un texte mais une suite de moutures du même
matériau. Les compter est utile, parce qu'elles ne disent pas la même chose et
que les fusionner reviendrait à effacer du contenu.

\begin{itemize}
\item \textbf{Page 24} — un premier jet, en langage de torseurs, qui pose
déjà tout : la donnée à l'ordre $2$, le groupe $G$, la conclusion
$G/G_0$. La phrase finale reste inachevée.
\item \textbf{Pages 25 et 26} — la mise en forme de théorème, avec ses items
a) et b), puis la correspondance biunivoque à trois termes. C'est la version
la plus articulée.
\item \textbf{Pages 27, 28 et 30} — le calcul explicite du changement de
paramètre. La page 30 est la plus complète des trois ; les pages 27 et 28 en
sont des états antérieurs.
\item \textbf{Page 29} — la preuve du théorème de structure de $G$, écrite
dans les termes géométriques.
\item \textbf{Page 31} — la même preuve, entièrement dépouillée : un lemme
sur un groupe $G \supset G_0$ opérant sur un ensemble $C \supset C_0$, sans
une ligne de géométrie.
\item \textbf{Pages 32 et 33} — le cadre reposé une dernière fois, avec la
définition de la régularité, et l'énoncé final.
\end{itemize}

\pagerange{1}{33}
\subsection*{Conventions, valables pour tout le dossier}

\textbf{Le groupe.} Les feuillets écrivent $\mathrm{PGL}(r)_S$ pour le groupe
structural d'un fibré de Brauer-Severi de dimension relative $r$. L'indice
compte donc la \emph{dimension relative} et non le rang du module : son
$\mathrm{PGL}(1)$ est le groupe des automorphismes de la droite projective,
que l'on écrit ici $H = \mathrm{PGL}_2$, de dimension $3$. La convention est
confirmée par le dossier lui-même : la formule d'Euler-Poincaré de la page 15
donne $\chi_{DR}(X, \mathfrak{g}) = 3\,\chi_{DR}(X, \mathcal{O}_X)$, ce qui
exige $\operatorname{rg} \mathfrak{g} = 3$.

\textbf{Le fibré et l'extension.} On suit la convention de EGA :
$P(E) = \operatorname{Proj} \operatorname{Sym}(E)$, de sorte qu'un point de
$P(E)$ au-dessus de $S'$ est un quotient inversible de $E$, et
$V(\Omega) = \operatorname{Spec} \operatorname{Sym}(\Omega)$, dont les
sections sont les homomorphismes $\Omega \to \mathcal{O}$, c'est-à-dire les
sections du \emph{dual} $\check{\Omega}$. C'est aussi ce que font les
pages 3 et 4, qui décrivent les points du complété formel par des
$\varphi : \Omega \to \mathcal{O}_S$ sans jamais nommer la
dualité.\footnote{La variance n'est pas un détail : c'est elle qui rend la
formule $\varphi^a = \varphi/(1+\varphi(a))$ de la page 4 cohérente avec
$u_a(\omega) = \omega/(1+a)$ de la même page. La première est une action sur
les points, la seconde sur les fonctions, et elles sont contragrédientes
l'une de l'autre.}

\textbf{Les modules.} $\mathcal{O}_S$, $\mathcal{O}_X$ pour ses $\mathcal{O}$
soulignés ; $\omega = \Omega^1_{X/S}$ sur une courbe relative ;
$\check{\Omega}$ pour le dual. $\Omega$ désigne, dans la partie formelle, le
module inversible $J/J^2$ ; il est identifié à $\Omega^1_{X/S}$ dès qu'une
connexion projective est en jeu, et il ne faut pas le confondre avec le
$\Omega = f_*(\omega)$ des pages 10 à 12, qui est de rang $g$. Les feuillets
emploient la même lettre pour les deux ; le contexte les sépare sans
ambiguïté, mais il faut le savoir.

\textbf{Les deux espaces de modules.} Les pages 14 et 15 écrivent
$M_g^{\mathfrak{z}}$ pour l'espace des couples (courbe de genre $g$, structure
projective), fibré sur $M_g$. Les pages 16 et 17 écrivent
$M_g^{\mathfrak{z}}(X)$ pour un objet tout différent : la variété des
caractères de $\pi = \pi_1(X,x)$, c'est-à-dire l'espace des représentations
$\pi \to H$ à conjugaison près. Les deux ont la même dimension $6g-6$, et
c'est le fait central ; ce ne sont pas le même espace, et
$M_g^{\mathfrak{z}}(X)$ n'est pas la fibre de $M_g^{\mathfrak{z}}$ au-dessus
de $X$ — celle-ci est de dimension $3g-3$. On garde ici les deux notations du
dossier, en réservant la parenthèse à la seconde.

\pagerange{14}{20}
\subsection*{Le glyphe de la structure projective}

À partir de la page 14 il note en exposant un signe que le transcripteur des
pages 1 à 20 n'a pas su nommer, et qu'il rend par $\mathfrak{z}$ en le
déclarant. C'est le même signe qui donne « $\mathfrak{z}$-structure » dans la
prose, et il y désigne la structure projective : $M_g^{\mathfrak{z}}$,
$\mathrm{Aut}^{\mathfrak{z}}(P)$, « $\mathfrak{z}$-fibré »,
« $\mathfrak{z}$-extension ».

Cette lecture reprend la convention telle quelle. Il faut dire clairement ce
qu'elle est : \textbf{$\mathfrak{z}$ est un choix de l'édition et non une
lecture du feuillet.} Personne n'a identifié le signe ; la lettre fraktur
n'est là que pour lui donner un corps typographique et pour qu'on puisse le
citer. Quiconque aura les feuillets sous les yeux pourra le nommer, et la
substitution sera alors mécanique.\footnote{Le transcripteur des pages 21 à 33
rencontre lui aussi, page 21, « un symbole isolé, qui revient plus bas devant
\emph{fibré principal} », et le laisse en blanc plutôt que de le rapporter au
$\mathfrak{z}$ du premier lot. Les deux observations sont indépendantes et
n'ont pas été rapprochées sur le fac-similé ; il est vraisemblable, mais non
établi, qu'il s'agisse du même signe.}

\pagerange{1}{33}
\subsection*{Ce que le contenu montre, et ce qu'il ne date pas}

Il faut le dire ici plutôt que de le taire. Le dossier appartient, dans
l'inventaire de Montpellier, au groupe 157-1 à 160, intitulé « Dérivateurs »
et daté 1990-[1991] ; la cote elle-même porte « s.d. ». Or rien sur ces
trente-trois pages ne relève des dérivateurs, ni des catégories dérivées, ni
de l'algèbre homotopique. Ce qui s'y trouve est du vocabulaire des années EGA
et SGA : fibrés de Brauer-Severi, $\mathrm{PGL}(r)_S$,
$\operatorname{Spf} \widehat{\operatorname{Sym}}$, parties principales
$P^n_{X/S}$ et $P^\infty_{X/S}$, connexion de Gauss-Manin, critère de descente
par les degrés. Les deux transcripteurs le notent séparément, chacun sur son
lot, sans s'être consultés.

Ce constat est solide en ceci qu'il est recontrôlable : il suffit de relire la
transcription. Il faut être exact sur sa portée, qui est étroite. \emph{Il ne
date rien.} Il dit de quoi les feuillets traitent, non quand ils ont été
écrits. Un mathématicien revient sur un sujet ancien quand il lui plaît, et un
dossier d'archives est une chemise de papier, groupée selon la place qu'elle
occupait et non selon ce qu'elle contient.

À côté de ce constat viennent des observations d'une autre espèce, et il faut
les en séparer nettement. Le transcripteur des pages 21 à 33 relève que les
pages 30, 32 et 33 sont écrites sur du papier à portées musicales imprimées ;
celui des pages 1 à 20 relève que d'autres feuillets sont du listing
informatique en accordéon, et que la page 4 porte une ligne tracée sur la
portée imprimée. Ce sont des observations matérielles, faites sur un
fac-similé numérique, où l'encre, la pression, le grain et le format ne se
jugent pas. Elles sont fragiles pour cette raison seule. Et quand bien même
elles seraient sûres, elles ne trancheraient rien : la récupération de papier
— portées, listings, versos — est une pratique qu'il a tenue sur des
décennies, et un support ne date pas ce qu'on écrit dessus.

En conséquence : \textbf{cette lecture ne conclut aucune date.} Le
La ligne de datation portée en tête est celle des archivistes, repris verbatim des
deux transcriptions, et il y reste. Ce qui est offert ici est la description
de ce que les feuillets montrent, pour que quelqu'un qui pourra les manipuler
sache quoi regarder.

\pagerange{1}{6}
\section*{Le dictionnaire : fibré projectif, groupe affine, extension (pages 1 à 6)}

Soit $S$ un schéma de base. Un \emph{fibré de Brauer-Severi} sur $S$ de
dimension relative $r$ est un schéma $P$ propre et lisse sur $S$ dont les
fibres géométriques sont des espaces projectifs de dimension $r$. Se donner un
tel $P$ revient à se donner un torseur sous $\mathrm{PGL}_{r+1}$ sur $S$ :
c'est la descente, pour la topologie fppf, du fait que
$\mathrm{Aut}(\mathbf{P}^r) = \mathrm{PGL}_{r+1}$.

\pagerange{1}{2}
\subsection*{Trois descriptions d'un même objet}

Se donner de plus une section $\sigma : S \to P$ revient à réduire le groupe
structural au stabilisateur d'un point de $\mathbf{P}^r$. Prenons pour point
celui qui correspond au quotient inversible $\mathcal{O}_S$ de
$\mathcal{O}_S^{\,r+1}$ défini par la dernière coordonnée ; son stabilisateur
est le groupe des matrices
\[
  \begin{pmatrix} A & C \\ 0 & 1 \end{pmatrix},
  \qquad A \in \mathrm{GL}_r, \quad C \in \mathbf{A}^r,
\]
c'est-à-dire de celles qui respectent la structure d'extension
$0 \to \mathcal{O}_S^{\,r} \to \mathcal{O}_S^{\,r+1} \to \mathcal{O}_S \to 0$
et induisent l'identité sur le quotient. Ce groupe est le \emph{groupe
affine}, produit semi-direct
\[
  \mathrm{Aff}(r) \;=\; \mathrm{GL}_r \ltimes \mathbf{G}_a^{\,r},
\]
de dimension $r^2 + r$ ; c'est bien $r$ de moins que
$\dim \mathrm{PGL}_{r+1} = r^2 + 2r$, comme il se doit puisque le quotient est
$\mathbf{P}^r$.\footnote{Le feuillet écrit $\mathrm{Aff}(r) =
\mathrm{GL}(r)\cdot E^r$ et précise « produit semi-direct » ; le point est de
ne pas lire ce produit comme direct, l'action de $\mathrm{GL}_r$ sur
$\mathbf{G}_a^r$ étant l'action tautologique. Dans la formule, l'argument de
$\mathrm{Aff}$ et celui de $\mathrm{GL}$ sont chacun écrits sur une lettre
biffée.}

Troisième description, celle dont tout le reste va se servir : la donnée de
$(P, \sigma)$ équivaut à celle d'une extension de modules localement libres
\[
  0 \longrightarrow \Omega \longrightarrow E
  \xrightarrow{\ p\ } \mathcal{O}_S \longrightarrow 0,
\]
avec $\Omega$ de rang $r$, le passage se faisant par $P = P(E)$ et
$\sigma$ la section définie par le quotient $p$. C'est une équivalence de
catégories, où les morphismes du côté des extensions sont les isomorphismes
induisant l'identité sur $\mathcal{O}_S$ — l'automorphisme d'une telle
extension est un couple (automorphisme de $\Omega$, élément de
$\mathrm{Hom}(\mathcal{O}_S, \Omega) = \Omega(S)$), ce qui redonne
exactement $\mathrm{Aff}(r)$.\footnote{La page 1 s'interrompt au milieu de
l'énoncé, qui est repris et achevé page 2, elle-même abandonnée et reprise du
même point de départ page 3. On donne ici l'énoncé complet, qui est celui des
pages 2 et 3 réunies. La précision sur les morphismes — que l'équivalence
porte sur la catégorie où $\Omega$ varie, et non sur les extensions d'un
$\Omega$ fixé — est ajoutée : sans elle l'équivalence serait fausse, car
$P(E) \simeq P(E \otimes L)$ pour tout $L$ inversible.}

\pagerange{3}{5}
\subsection*{Le complété formel le long de la section, et l'action de $\Omega$}

Soient $\widehat{P}$ le complété formel de $P$ le long de $\sigma(S)$,
$\widehat{S} = \mathcal{O}_{\widehat{P}}$ son algèbre, $\widehat{J}$ l'idéal
d'augmentation, et $L = \mathcal{O}_P(1)$ avec $\widehat{L}$ sa restriction.
Le groupe $\Omega$, vu comme
$\mathrm{Hom}(\mathcal{O}_S, \Omega) \subset \mathrm{Aff}(r)$, opère sur
l'extension, donc sur $(P,\sigma)$, donc sur $\widehat{S}$ et sur
$\widehat{L}$, de façon compatible.

Supposons l'extension scindée, $E \simeq \Omega \oplus \mathcal{O}_S$. Alors
les points de $\widehat{P}$ à valeurs dans $S'$ sont les quotients inversibles
$L'$ de $E$ qui coïncident, sur le réduit, avec le quotient donné ; en
normalisant $L'$ par l'image de la section $e$ de $\mathcal{O}_S$, un tel
quotient s'écrit $(x,\lambda) \mapsto \varphi(x) + \lambda$ avec
$\varphi : \Omega \to \mathcal{O}_{S'}$ à valeurs dans le nilradical. D'où
l'identification
\[
  \widehat{P} \;\simeq\; V(\Omega)^{\wedge}
  \;=\; \operatorname{Spf} \widehat{\operatorname{Sym}}(\Omega),
\]
le complété de $V(\Omega)$ le long de la section nulle.

L'action de $a \in \Omega(S)$ est alors entièrement calculable. Sur
l'extension, $u_a(x,\lambda) = (x + \lambda a, \lambda)$. Sur les points,
\[
  \varphi^a \;=\; u_a^*(\varphi) \;=\; \frac{\varphi}{1 + \varphi(a)} .
\]
Sur les fonctions, $u_a$ est l'automorphisme d'algèbre de
$\widehat{\operatorname{Sym}}(\Omega)$ qui envoie $\omega \in \Omega$ sur
\[
  u_a(\omega) \;=\; \frac{\omega}{1+a}
  \;=\; \sum_{n \geqslant 0} (-1)^n\, \omega\, a^n ,
\]
le quotient ayant un sens parce que $1 + a$ est inversible dans le complété.
Comme $u_a$ est un morphisme d'algèbres, on en déduit la règle de degré, que
le feuillet encadre :
\[
  u_a(f) \;=\; \frac{f}{(1+a)^n} \qquad \text{si } \deg f = n .
\]
Sur le module inversible $\widehat{L}$ la règle est décalée d'un cran. On
identifie $\widehat{L}$ à $\widehat{S}$ par l'épimorphisme canonique
$\pi : E \otimes_{\mathcal{O}_S} \widehat{S} \to \widehat{L}$, qui s'écrit
$\pi(\omega \otimes f + e \otimes g) = \omega f + g$ ; la commutation demandée
$\pi \circ u_{a*}^{\widehat{E}} = u_{a*}^{\widehat{L}} \circ \pi$ force alors
\[
  u_{a*}^{\widehat{L}}(g) \;=\; (1+a)\, u_{a*}(g)
  \;=\; \frac{g}{(1+a)^{\,n-1}} \qquad \text{si } \deg g = n .
\]
Le feuillet vérifie l'identité sur un élément général, et le calcul tombe
juste.\footnote{Dans le cas général d'une extension quelconque, définie par un
torseur $T$ sous $\Omega$, on obtient $\widehat{S}$ et $\widehat{L}$ comme
produits contractés de $T$ par $\widehat{\operatorname{Sym}}(\Omega)$ le long
des deux actions ci-dessus. Le symbole de contraction est douteux sur le
feuillet et les dernières lignes de la page 5 sont biffées ; on n'en tire rien
de plus que le principe.}

\pagerange{6}{6}
\subsection*{La conséquence gradue}

La page 6 est un brouillon au crayon très pâle, traversé par la transparence
du verso : les formules portent, la prose qui les relie est en grande partie
perdue. Ce qui s'en tire est ceci. L'idéal d'augmentation $\widehat{J}$ donne
$\operatorname{gr}^1_{\widehat{J}}(\widehat{S}) \simeq \Omega$, donc
$\operatorname{gr}_{\widehat{J}}(\widehat{S}) \simeq
\operatorname{Sym}(\Omega)$ et $J^n/J^{n+1} \simeq \operatorname{Sym}^n\Omega$
canoniquement. Ce qui n'est \emph{pas} canonique, et qui porte l'information,
est le cran suivant : $J^n/J^{n+2}$ est une extension de
$\operatorname{Sym}^n\Omega$ par $\operatorname{Sym}^{n+1}\Omega$, et sa
classe s'obtient à partir du torseur $T$ par l'homomorphisme
\[
  \Omega \longrightarrow
  \mathrm{Hom}\big( \operatorname{Sym}^n \Omega,\;
  \operatorname{Sym}^{n+1} \Omega \big),
  \qquad a \longmapsto \big( f \mapsto -\,n\,a f \big).
\]
Le facteur $n$ et le signe sont ce que produit la règle de degré :
$u_a(f) = f(1+a)^{-n} \equiv f - n\,a f$ modulo le degré
$n+2$.\footnote{Le feuillet écrit la flèche deux fois et pas de la même
manière : en haut $f \mapsto n a f + \ldots$, plus bas $v(a)(f) = af$ sans
coefficient. C'est la page 7 qui tranche : elle nomme la flèche
$-u_1$, c'est-à-dire la valeur en $n=1$ de l'expression ci-dessus, signe
compris. On a donc retenu $a \mapsto (f \mapsto -naf)$, qui est la seule des
trois écritures compatible avec le reste du dossier. Le reste de la page 6 —
une demi-douzaine de membres de phrase — est perdu.}

\pagerange{7}{9}
\section*{La définition d'une connexion projective (pages 7 à 9)}

\pagerange{7}{7}
\subsection*{Les deux données, et l'identification qu'elles forcent}

Soit $X$ un schéma lisse sur $S$. Une \emph{connexion projective} sur $X$
relativement à $S$ est la donnée de :

\begin{itemize}
\item[a)] un fibré de Brauer-Severi $P$ sur $X$ muni d'une section,
c'est-à-dire, par le dictionnaire des pages 1 à 6, une extension de modules
localement libres sur $X$
\[
  0 \longrightarrow \Omega \longrightarrow E \longrightarrow
  \mathcal{O}_X \longrightarrow 0 \,;
\]
\item[b)] pour chaque $n$, un isomorphisme de $\mathcal{O}_X$-algèbres
augmentées
\[
  P^n_{X/S} \;\simeq\; \mathcal{O}_{\widehat{P}} \big/ \widehat{J}^{\,n+1},
\]
compatible au passage de $n+1$ à $n$, où $P^n_{X/S}$ est le faisceau des
parties principales d'ordre $n$ — l'algèbre du $n$-ième voisinage
infinitésimal de la diagonale de $X \times_S X$ — et $\widehat{P}$ le complété
formel de $P$ le long de la section.
\end{itemize}

C'est la traduction algébrique littérale de « une carte à valeurs dans
$\mathbf{P}^r$ » : la donnée b) dit que le voisinage infinitésimal d'un point
de $X$ est identifié, à tout ordre fini, au voisinage infinitésimal d'un point
de $\mathbf{P}^r$, et que ces identifications sont cohérentes entre elles. Il
n'y a pas de convergence à demander : c'est tout l'intérêt de la version
algébrique.

La donnée b) à l'ordre $1$ ne coûte rien mais impose tout : les gradués de
degré $1$ des deux membres sont $\Omega^1_{X/S}$ d'un côté et $\Omega$ de
l'autre, d'où l'identification
\[
  (*) \qquad \Omega \;\simeq\; \Omega^1_{X/S} .
\]
En particulier $r = \dim X/S$, et $E$ est une extension de $\mathcal{O}_X$ par
$\Omega^1_{X/S}$.

Le premier contenu réel est à l'ordre $2$. Du côté des parties principales,
l'idéal d'augmentation donne une extension
\[
  0 \longrightarrow \operatorname{Sym}^2 \Omega^1_{X/S} \longrightarrow
  I/I^3 \longrightarrow \Omega^1_{X/S} \longrightarrow 0
\]
dont la classe vit dans
$\mathrm{Ext}^1(\Omega^1_{X/S}, \operatorname{Sym}^2\Omega^1_{X/S})$, et un
isomorphisme comme en b) est en particulier un scindage relatif, c'est-à-dire
un point d'un torseur sous
$\mathrm{Hom}(\Omega^1_{X/S}, \operatorname{Sym}^2 \Omega^1_{X/S})$. Du côté
de $\widehat{P}$, la page 6 vient de calculer l'extension analogue : c'est
celle déduite du torseur $T$ par $-u_1$. La condition b) à l'ordre $2$ dit
donc que les deux classes se correspondent, et fixe de plus
l'isomorphisme.\footnote{La page 7 est au crayon, pâle, avec les dernières
lignes reprises à l'encre bleue ; elle porte la charge d'illisible la plus
lourde du premier lot avec les pages 6, 11 et 21. Les mots qui articulent la
phrase — « est déterminée par », « d'un cocycle », « doit être muni d'un
iso » — sont là, mais ce qui les relie ne l'est pas. L'énoncé donné ci-dessus
est la seule lecture cohérente de ce qui reste ; il n'est pas une
reconstruction du feuillet, et on ne comble pas les blancs.}

\pagerange{8}{9}
\subsection*{Intégrabilité, et l'écriture en coordonnées}

L'algèbre $P^\infty_{X/S} = \varprojlim_n P^n_{X/S}$ porte une
comultiplication canonique
\[
  \Delta : P^\infty_{X/S} \longrightarrow
  P^\infty_{X/S} \otimes_{\mathcal{O}_X} P^\infty_{X/S}
\]
et une augmentation vers $\mathcal{O}_X$ : c'est le groupoïde formel du
voisinage de la diagonale. On dit qu'une connexion projective à l'ordre
infini est \emph{intégrable} si elle est compatible à cette structure,
autrement dit si elle provient d'une connexion sur $P = P(E)$ au sens
ordinaire ; cette dernière est alors déterminée de façon
unique.\footnote{Le feuillet pose l'intégrabilité exactement ainsi — « si elle
[provient] d'une $\infty$-connexion de $P = P(E)$, laquelle sera dès lors
déterminée » — mais les verbes sont illisibles. On a choisi la formulation
minimale. C'est la même page qui pose, sous forme de question et sans y
répondre, celle que le titre du dossier annonce : une telle connexion est-elle
canoniquement déterminée par $X/S$ ? Voir la section finale.}

La condition étant locale sur $X$, on se place au-dessus d'un ouvert muni de
coordonnées locales $f_1, \ldots, f_n$ étales sur $S$. Alors
\[
  P^\infty_{X/S} \;\simeq\;
  \mathcal{O}_X[[\delta f_1, \ldots, \delta f_n]],
  \qquad \delta f_i = f_i \otimes 1 - 1 \otimes f_i,
\]
et la comultiplication est primitive sur les générateurs :
\[
  \Delta(\delta f_i) = \delta f_i \otimes 1 + 1 \otimes \delta f_i .
\]
Symétriquement, $\Omega \simeq \Omega^1_{X/S} \simeq \mathcal{O}_X^{\,n}$
donne $\widehat{S} \simeq \mathcal{O}_X[[T_1, \ldots, T_n]]$.\footnote{Le
feuillet écrit $\mathcal{O}_S[[\delta f_1, \ldots]]$ et
$\mathcal{O}_S[[T_1, \ldots, T_n]]$. C'est une inadvertance : $P^\infty_{X/S}$
et $\widehat{S}$ sont des $\mathcal{O}_X$-algèbres, les $\delta f_i$ et les
$T_i$ engendrent au-dessus de $\mathcal{O}_X$, et l'algèbre $\mathcal{O}_S$ ne
donnerait pas les bons gradués. On a corrigé.}

La connexion est alors la donnée des images réciproques des générateurs :
\[
  F_i = \delta f_i + G_i \in \Gamma\, P^\infty_{X/S},
  \qquad \Phi_i = T_i + \Psi_i \in \Gamma\, \widehat{S},
\]
avec $G_i$ et $\Psi_i$ d'ordre $\geqslant 2$, et l'intégrabilité est la
compatibilité de $c_\infty : P^\infty_{X/S} \to \widehat{S}$ aux deux
comultiplications. Le feuillet s'arrête au premier tiers, sur le mot
« Posant », au moment précis d'écrire cette compatibilité ; la suite manque et
on ne la supplée pas.

\pagerange{10}{12}
\section*{Dimension relative un : de Rham, autodualité, Kodaira-Spencer (pages 10 à 12)}

\pagerange{10}{10}
\subsection*{La cohomologie de de Rham d'une famille de courbes}

Soit maintenant $f : X \to S$ propre et lisse de dimension relative $1$, à
fibres connexes de genre $g$. On pose
\[
  \omega = \Omega^1_{X/S}, \qquad \Omega = f_*(\omega), \qquad
  DR = R^1 f_*\big( \Omega^\bullet_{X/S} \big),
\]
de sorte que $\Omega$ est localement libre de rang $g$ et $DR$ de rang $2g$.
La dualité de Serre donne $R^1 f_*(\mathcal{O}_X) \simeq \check{\Omega}$, et
la suite spectrale de Hodge vers de Rham se réduit à la suite exacte
\[
  0 \longrightarrow \Omega \longrightarrow DR \longrightarrow
  \check{\Omega} \longrightarrow 0 .
\]
Le cup-produit $\Lambda^2 DR \to R^2 f_*(\Omega^\bullet_{X/S}) \simeq
\mathcal{O}_S$ est un accouplement alterné parfait, pour lequel $\Omega$ est
un sous-module lagrangien : c'est ce que la marge du feuillet résume d'un mot,
« autodual ».

Les scindages de la suite de Hodge forment un torseur sous
$\mathrm{Hom}(\check{\Omega}, \Omega) = \Omega \otimes \Omega$. Un scindage
$q : \check{\Omega} \to DR$ donne $DR \simeq \Omega \oplus \check{\Omega}$, et
dans cette décomposition la forme symplectique s'écrit matriciellement
\[
  \begin{pmatrix} 0 & \pm\,\mathrm{id} \\ \mathrm{id} & \varphi_0 \end{pmatrix},
\]
où $\varphi_0$ est la forme alternée induite sur $\check{\Omega}$. Un autre
choix $q' = q + u$ la change selon
\[
  \varphi'_0(t,t') = \varphi_0(t,t') + \langle t, u(t')\rangle
  - \langle t', u(t)\rangle ,
\]
c'est-à-dire du double de l'antisymétrisé de $u$. Il en résulte que les
scindages \emph{lagrangiens} — ceux pour lesquels $\varphi_0 = 0$ — existent
dès que $2$ est inversible, et qu'ils forment un torseur sous la partie
symétrique $\operatorname{Sym}^2 \Omega$, noyau de
$\Omega \otimes \Omega \to \Lambda^2 \Omega$.\footnote{C'est la suite exacte
que la page 12 écrit sous la forme
$0 \to \Gamma^2\Omega \to \Omega\otimes\Omega \to \Lambda^2\Omega \to 0$, avec
la note marginale « Epi (car. $\neq 2$) ». $\Gamma^2$ est le carré à
puissances divisées, c'est-à-dire le module des tenseurs symétriques ; en
caractéristique différente de $2$ il s'identifie à $\operatorname{Sym}^2$, et
c'est là toute la portée de la restriction sur la caractéristique. Le signe
entre les deux derniers termes de la formule de variation est masqué par un
pâté d'encre sur le feuillet ; le transcripteur lit $-$, et c'est le signe que
l'antisymétrie demande.}

\pagerange{11}{11}
\subsection*{Gauss-Manin, et ce que la page 11 ne livre pas}

La page 11 est à l'encre bleue, d'une écriture rapide : les formules portent,
la prose qui les relie est en grande partie perdue. C'est, avec les pages 6, 7
et 21, l'un des quatre feuillets les plus abîmés du dossier. Voici ce qui s'en
tire sans reconstruction.

Si $S$ est lui-même relatif à une base $T$, la connexion de Gauss-Manin
\[
  \nabla : DR \longrightarrow DR \otimes \Omega^1_{S/T}
\]
n'est pas compatible à la filtration de Hodge, et ses défauts de compatibilité
sont $\mathcal{O}_S$-linéaires : le premier est un homomorphisme
$\Omega \to \check{\Omega} \otimes \Omega^1_{S/T}$, soit encore, par
adjonction, un homomorphisme depuis ou vers $\check{\Omega} \otimes
\check{\Omega}$. C'est la transversalité de Griffiths, et l'application
obtenue est celle de Kodaira-Spencer.\footnote{Les flèches de la page 11 sont
écrites $\Omega^1_{S/T} \to \check{\Omega}\otimes\check{\Omega}$ et
$\Omega^1_{S/T} \to \Lambda^2\check{\Omega}$. Telles quelles, la source et le
but sont dans la variance qu'aurait l'application duale : c'est
$T_{S/T} \to \operatorname{Sym}^2\check{\Omega}$ qui est l'application de
Kodaira-Spencer, et $\Omega^1_{S/T} \simeq f_*(\omega^{\otimes 2})$ qui en est
la transposée — c'est d'ailleurs ce que la page 12 écrit. Le feuillet est trop
abîmé pour qu'on décide s'il écrit l'application ou sa duale, et cette édition
ne tranche pas. De même, la seconde flèche a pour but $\Lambda^2$ là où
l'application de Kodaira-Spencer est symétrique ; on ne bâtit rien dessus.}

Le reste de la page — une remarque sur les torseurs sous un groupe hyperbolique
et le sous-groupe $\pi_1$ des automorphismes, une suite exacte faisant
intervenir $H^0(X, \Omega_X)$ et un sous-fibré de corang $1$ sur une courbe de
genre $\geqslant 2$ — n'est pas lisible assez continûment pour qu'on en tire
un énoncé. On le signale et on passe.

\pagerange{12}{12}
\subsection*{Le compte des quadriques}

La page 12 n'est pas de la prose : c'est une colonne de formules jetées. Deux
d'entre elles portent, et elles se répondent.

La première est
\[
  f_*(\omega^{\otimes 2}) \;\simeq\; \Omega^1_{S/T} .
\]
C'est l'isomorphisme de Kodaira-Spencer pour les modules de courbes : il dit
que $S \to T$ est versel, c'est-à-dire que $S$ est (localement) l'espace de
modules $M_g$. Comme $\deg \omega^{\otimes 2} = 4g-4$ et
$H^1(X, \omega^{\otimes 2}) = 0$ pour $g \geqslant 2$, Riemann-Roch donne
$\operatorname{rg} f_*(\omega^{\otimes 2}) = 3g-3$ : c'est la dimension de
$M_g$.

La seconde est le compte
\[
  \frac{g(g+1)}{2} - (3g-3) \;=\; \frac{1}{2}\big(g^2 - 5g + 6\big)
  \;=\; \frac{1}{2}(g-2)(g-3) .
\]
Les deux termes de gauche sont les rangs de
$\operatorname{Sym}^2 \Omega = \operatorname{Sym}^2 f_*(\omega)$ et de
$f_*(\omega^{\otimes 2})$, et la flèche qui va du premier au second est la
multiplication des sections. Le compte est donc la dimension de son noyau
lorsqu'elle est surjective, c'est-à-dire — c'est le théorème de Max Noether —
lorsque $X$ n'est pas hyperelliptique. Et ce noyau a un nom classique : c'est
l'espace des quadriques de $\mathbf{P}^{g-1}$ passant par la courbe
canonique, de dimension $\binom{g-2}{2}$.\footnote{L'identification des deux
termes du compte est le fait de cette édition. Le feuillet n'écrit que
l'arithmétique ; ce qui la justifie est son voisinage immédiat — la suite
$0 \to \Gamma^2\Omega \to \Omega\otimes\Omega \to \Lambda^2\Omega \to 0$ deux
lignes plus haut, et $f_*(\omega^2) \simeq \Omega^1_{S/T}$ une ligne plus
bas — et rien d'autre. La mention de Max Noether est nôtre ; le feuillet ne le
nomme pas, et ne dit pas non plus s'il suppose $X$ non hyperelliptique.}

L'intérêt du compte pour le dossier est le suivant. Les scindages lagrangiens
de $DR$ forment un torseur sous $\operatorname{Sym}^2\Omega$ ; les connexions
projectives, on va le voir, forment un torseur sous
$f_*(\omega^{\otimes 2})$ ; et la flèche de l'un vers l'autre est la
multiplication. Le compte mesure donc exactement de combien un scindage de
Hodge détermine \emph{plus} qu'une connexion projective. Le reste de la
page — la suite $0 \to \mathcal{O}_X \to P^1 \to \omega \to 0$ des parties
principales, le nombre $2-2g$, les modules $f_*(\omega^n \otimes P)$ et leur
gradué — est du repérage sans énoncé.

\pagerange{14}{18}
\section*{L'espace de modules des structures projectives (pages 14 à 18)}

\pagerange{14}{14}
\subsection*{L'espace, et le critère de non-dégénérescence}

On fixe un genre $g$ et un foncteur rigidifiant sur les courbes propres et
lisses relatives de genre $g$ — de Torelli, ou de niveau convenable — de
manière que $M_g$ soit un espace et non seulement un champ. Pour $S$ de
caractéristique $0$, on pose alors : $M_g^{\mathfrak{z}}(S)$ est l'ensemble
des courbes propres lisses $X/S$ de genre $g$, munies d'un fibré en droites
projectives $P$ sur $X$ à structure projective, et soumises à la seule
condition $\mathrm{Aut}^{\mathfrak{z}}(P_s) = 1$ pour tout point $s$ de
$S$.\footnote{Le glyphe $\mathfrak{z}$ est
celui de l'édition : voir la section qui lui est consacrée. La condition
fibre à fibre est ce que le feuillet appelle avoir un « point
$\mathfrak{z}$-ordinaire », en opposition à « dégénéré » ; le sens en est
fixé par la remarque qui suit immédiatement.}

Le feuillet explicite aussitôt ce que la condition exclut. Soient $X$ un
schéma algébrique connexe sur $k$ algébriquement clos et $P$ un fibré
projectif à structure projective sur $X$. Considérons :

\begin{itemize}
\item[a)] $\mathrm{Aut}^{\mathfrak{z}}(P) = 1$ ;
\item[b)] le groupe algébrique
$\mathrm{Aut}^{\mathfrak{z}}_{X/k}(P) \subset H = \mathrm{PGL}_2$ est de
dimension $0$ ;
\item[c)] $P$ n'est d'aucun des deux types dégénérés ci-dessous.
\end{itemize}

Les deux types dégénérés sont :

\begin{itemize}
\item[$\alpha$)] le type $\mathbf{G}_m$ : $P \simeq P(\Omega \oplus
\mathcal{O}_X)$, clôture projective du fibré en droites $V(\Omega)$, avec sa
structure évidente ;
\item[$\beta$)] le type $\mathbf{G}_a$ : $P \simeq P(E)$ pour une extension
$0 \to \mathcal{O}_X \to E \to \mathcal{O}_X \to 0$ à structure projective,
c'est-à-dire la clôture projective d'un torseur sous $\mathbf{G}_a$.
\end{itemize}

Ce qui est établi est b) $\Leftrightarrow$ c) — un sous-groupe algébrique de
$\mathrm{PGL}_2$ de dimension $> 0$ contient un $\mathbf{G}_m$ ou un
$\mathbf{G}_a$, et le fibré est alors la clôture projective de la
représentation correspondante — et a) $\Rightarrow$ b), trivialement. Le
feuillet donne les trois pour équivalentes ; telles qu'elles sont écrites,
a) est strictement plus forte que b), car un groupe fini non trivial est de
dimension $0$.\footnote{L'exemple est immédiat : le groupe de Klein
$V_4 \subset \mathrm{PGL}_2$ est fini, non trivial, et son propre
centralisateur. Un fibré projectif plat de monodromie $V_4$ vérifie b) et c)
sans vérifier a). Cette lecture énonce donc b) $\Leftrightarrow$ c) et
a) $\Rightarrow$ b), et laisse l'équivalence complète à la condition
supplémentaire qui la rend vraie — voir la remarque ci-dessous. Le feuillet
n'énonce aucune restriction sur $X$ à cet endroit, ce qui est précisément d'où
vient l'écart.}

\textbf{Remarque de l'édition.} Dans le seul cas dont le dossier ait
finalement besoin — $X$ courbe projective lisse de genre $\geqslant 2$ sur un
corps de caractéristique $0$, munie d'une connexion projective — les trois
conditions \emph{sont} équivalentes, et c'est le théorème des pages 21 et 22
qui le donne : il y est montré que le groupe structural ne se réduit à aucun
sous-groupe propre, donc que l'adhérence de Zariski de la monodromie est $H$
tout entier, donc que le centralisateur de la monodromie est le centre de
$\mathrm{PGL}_2$, qui est trivial. Le dossier ne fait jamais ce rapprochement ;
il est de cette édition, et il est forcé par les deux énoncés tels qu'ils
sont.

\pagerange{15}{15}
\subsection*{Le calcul de dimension}

Soit $\mathfrak{g} = \mathrm{ad}(P)$ le fibré en algèbres de Lie adjoint,
localement libre de rang $3$, muni de sa connexion plate. Le feuillet propose
trois conditions supplémentaires, pour $X$ courbe propre lisse de genre $g$ :

\begin{itemize}
\item[d)] $\dim H^1_{DR}(X, \mathfrak{g}) = 6g-6$ ;
\item[d$'$)] $H^0_{DR}(X, \mathfrak{g}) = 0$ ;
\item[d$''$)] $H^2_{DR}(X, \mathfrak{g}) = 0$.
\end{itemize}

Leur équivalence tient en deux lignes. La forme de Killing donne
$\mathfrak{g} \simeq \check{\mathfrak{g}}$, d'où par dualité de Poincaré
\[
  (*) \qquad H^0_{DR}(X, \mathfrak{g}) \;\simeq\;
  H^2_{DR}(X, \mathfrak{g})^{\vee},
\]
ce qui donne d$'$) $\Leftrightarrow$ d$''$). La caractéristique
d'Euler-Poincaré de de Rham ne dépend que du rang :
\[
  (**) \qquad \chi_{DR}(X, \mathfrak{g}) = 3\,\chi_{DR}(X, \mathcal{O}_X)
  = 3(2-2g) = 6-6g .
\]
Combinant, on obtient
\[
  (\!*\!*\!*) \qquad \dim H^1_{DR}(X, \mathfrak{g})
  = 2 \dim H^0_{DR}(X, \mathfrak{g}) + (6g-6),
\]
d'où d) $\Leftrightarrow$ d$'$). Enfin $H^0_{DR}(X, \mathfrak{g}) \simeq
\operatorname{Lie} \mathrm{Aut}^{\mathfrak{z}}_{X/k}(P)$ — une section
horizontale du fibré adjoint est un automorphisme infinitésimal — d'où
d$'$) $\Leftrightarrow$ b), et la formule
\[
  \dim H^1_{DR}(X, \mathfrak{g})
  = 2\Big( \dim \mathrm{Aut}^{\mathfrak{z}}_{X/k}(P) + 3g-3 \Big).
\]
Le feuillet porte, sous le membre de gauche, l'accolade qui dit ce qu'est cet
espace : « espace tangent à l'espace modulaire $M_g^{\mathfrak{z}}$ au point
correspondant à $(X, P, \mathfrak{z})$ ».

De là : la condition d$'$) signifie que $M_g^{\mathfrak{z}}$ n'a pas
d'automorphismes infinitésimaux et c) qu'il n'a pas d'automorphismes finis, de
sorte que $M_g^{\mathfrak{z}}$ est rigide — et donc représentable, là où $M_g$
ne l'est qu'après rigidification. La condition d$''$) donne la lissité
formelle de $M_g^{\mathfrak{z}}$ au-dessus de $M_g$. Et la condition d) donne
la dimension.

\pagerange{15}{15}
\subsection*{L'anomalie de la page 15, et ce que la cohérence demande}

Le feuillet écrit : « d$''$) implique que $M_g^{\mathfrak{z}}$ est
formellement lisse sur $M_g$, et d) que sa dimension relative sur $M_g$ est
$6g-6$ ». Le transcripteur note expressément que la dimension relative se lit
bien $6g-6$, et il l'a transcrite telle quelle, sans réparation. Or la
formule qui précède immédiatement donne $3g-3$.

Cette lecture ne corrige pas le feuillet. Elle dit ce que la cohérence
demande, et sur quoi elle se fonde.

Ce que la cohérence demande :
\[
  \dim M_g^{\mathfrak{z}} = 6g-6, \qquad \dim M_g = 3g-3, \qquad
  \dim \big( M_g^{\mathfrak{z}} / M_g \big) = 3g-3 .
\]

Ce sur quoi elle se fonde — quatre appuis, tous internes au dossier :

\begin{itemize}
\item L'accolade de la page 15 elle-même dit que $H^1_{DR}(X,\mathfrak{g})$
est l'espace tangent à $M_g^{\mathfrak{z}}$, sans mention de $M_g$. Le
$6g-6$ est donc la dimension \emph{absolue}.
\item La page 12 donne $f_*(\omega^{\otimes 2}) \simeq \Omega^1_{S/T}$, de
rang $3g-3$ : c'est la dimension de $M_g$, et le dossier s'en sert dans le
compte $g(g+1)/2 - (3g-3)$ de la même page.
\item La fibre de $M_g^{\mathfrak{z}} \to M_g$ au-dessus d'une courbe $X$ est
l'ensemble des connexions projectives sur $X$, dont les pages 26, 29, 31 et 33
établissent qu'il est un torseur sous $\Gamma(X, \Omega^{\otimes 2}) =
H^0(X, \omega_X^{\otimes 2})$, de dimension $3g-3$ pour $g \geqslant 2$. C'est
le résultat principal du dernier tiers du dossier.
\item La page 17 nomme le même module, $H^0(X, \omega^{\otimes 2}_{X/\mathbf{C}})$,
comme celui qui gouverne l'espace $M_{\mathrm{ann}}(X)$ des structures
projectives sur une courbe analytique fixée ; et la page 20 en donne l'action
explicite, $u_\omega(c) = c + d^1_{\omega/S}(\omega)$.
\end{itemize}

Autrement dit : $6g-6 = (3g-3) + (3g-3)$, base plus fibre, et le dossier
calcule les trois nombres. L'explication la plus économique de ce qu'on lit
page 15 est qu'il a écrit la dimension absolue dans une phrase qui parlait de
la dimension relative ; une autre est que les mots « relative sur $M_g$ » ne
portaient, dans son intention, que sur la lissité. Cette édition n'en choisit
aucune et ne réécrit pas le feuillet : ce qui est sur le papier est $6g-6$, et
la transcription le dit.

\pagerange{16}{17}
\subsection*{Le tableau analytique : monodromie, Teichmüller, uniformisation}

On passe sur $\mathbf{C}$. Soient $X$ une courbe analytique compacte connexe
de genre $g \geqslant 2$, $x \in X$, $\pi = \pi_1(X,x)$ et
$H = \mathrm{PGL}_2(\mathbf{C})$. Le feuillet pose
\[
  M_g^{\mathfrak{z}}(X) \;\simeq\;
  \mathrm{Hom}^{\,!}(\pi, H) \,\big/\, H,
\]
où $\mathrm{Hom}^{\,!}$ désigne l'ensemble des homomorphismes dont le
stabilisateur dans $H$ — c'est-à-dire le centralisateur de l'image — se réduit
à l'élément neutre. C'est la variété des caractères, lisse de dimension
$(2g-2)\dim H = 6g-6$ : le nombre coïncide avec celui de la page 15 sans que
les deux espaces coïncident, et c'est tout l'intérêt de la
chose.\footnote{C'est la correspondance de Riemann-Hilbert. La page 16 décrit
$M_g^{\mathfrak{z}}(X)$ comme « l'espace analytique associé à une solution sur
$\mathbf{C}$, lisse de dimension $6g-6$ », et laisse implicite qu'il ne dépend
que de $\pi$, donc ni de la structure complexe de $X$, ni du fibré. La
condition « $\,!$ » est ce qui fait du quotient un espace, et c'est
exactement la condition a) de la page 14, transportée par la monodromie.}

Le groupe $\mathrm{Aut}(\pi)$ opère à droite sur $\mathrm{Hom}(\pi,H)$ par
$\rho \cdot u = \rho \circ u$ ; comme $\rho \circ \mathrm{int}(\gamma) =
\mathrm{int}(\rho(\gamma)) \circ \rho$, les automorphismes intérieurs
agissent trivialement sur le quotient, et c'est donc
\[
  \mathrm{Out}(\pi) = \mathrm{Aut}(\pi)/\mathrm{Int}(\pi)
\]
— le groupe modulaire de la surface — qui opère sur
$M_g^{\mathfrak{z}}(X)$.

Vient alors la pièce qui justifie le mot « uniformisation » du titre. Si $Y$
est une courbe analytique complète de genre $g$ munie d'un isomorphisme
$\pi_1(Y) \simeq \pi$, l'uniformisation de $Y$ par le disque fournit une
représentation $\rho_Y : \pi \to H$, définie à conjugaison près, donc un point
de $M_g^{\mathfrak{z}}(X)$. On obtient ainsi un morphisme analytique
\emph{injectif}
\[
  M_g(X) \hookrightarrow M_g^{\mathfrak{z}}(X),
\]
où $M_g(X)$ est l'espace de Teichmüller rigidifié par $\pi$, et ce morphisme
commute aux actions du groupe modulaire $\Gamma$.

La colonne de droite du même feuillet caractérise l'image, et c'est là que le
dossier touche l'uniformisation de plus près : ce sont les
$\rho : \pi \hookrightarrow H$ injectifs pour lesquels il existe un domaine
$D \subset \mathbf{P}^1(\mathbf{C})$ invariant — unique s'il existe — sur
lequel $\pi$ opère librement avec quotient compact. C'est la définition même
d'une uniformisation, réécrite dans la variété des caractères. Mais le
feuillet ouvre la phrase par « on suppose » : rien n'est démontré ici, ni la
constructibilité de l'image, ni son caractère algébrique.\footnote{La colonne
de droite de la page 17 est lourdement raturée et sa moitié inférieure,
surchargée de renvois, ne se laisse pas suivre. Les mots « on suppose que
l'image [est] constructible [et] algébrique » sont ce qu'on lit, et ils sont
présentés comme une supposition. La lecture ne va pas au-delà.}

Le feuillet définit enfin $M_{\mathrm{ann}}(X)$, l'espace des structures
projectives sur la courbe analytique fixée $X$, et note que
$H^0(X, \omega^{\otimes 2}_{X/\mathbf{C}})$ le gouverne ; puis il pose le
morphisme
\[
  M_{\mathrm{ann}}(X) \longrightarrow M_g^{\mathfrak{z}}(X)
\]
et demande : « est-il injectif ? » C'est l'application de monodromie des
structures projectives, dont le comportement local et l'injectivité sont
devenus un sujet en soi dans la littérature ; le feuillet pose la question et
ne la tranche pas, et cette édition non plus.\footnote{La phrase qui décrit
$M_{\mathrm{ann}}(X)$ est mutilée : on y lit « c'est un [\ldots] sur le
[\ldots] de deux fois $H^0(X,\omega^{\otimes 2}_{X/\mathbf{C}})$ », et ni le
mot que « deux fois » qualifie, ni la nature de la structure — torseur,
fibré — ne se laissent lire. On n'a retenu que le module, qui est certain
parce qu'il est écrit en toutes lettres.}

\pagerange{18}{18}
\subsection*{Les automorphismes du couple $(P,\sigma)$}

Le tiers supérieur de la page 18 est annulé par des hachures. Ce qui suit se
tient. Soit $E$ l'extension non scindée $0 \to \omega \to E \to \mathcal{O}_X
\to 0$, de classe $c \in \mathrm{Ext}^1(\mathcal{O}_X, \omega) =
H^1(X, \omega) \simeq k$, avec $c \neq 0$. Un automorphisme $u$ de $E$
préservant $\omega$ agit sur $\omega$ par un scalaire $\lambda \in
H^0(X, \mathcal{O}_X^*) = \mathbf{C}^*$, et multiplie donc la classe $c$ par
$\lambda$ ; de $\lambda c = c$ et $c \neq 0$ on tire $\lambda = 1$. Un tel $u$
est donc un automorphisme d'extension, c'est-à-dire un homomorphisme
$\mathcal{O}_X \to \omega$, c'est-à-dire une forme différentielle globale.
D'où
\[
  \mathrm{Aut}_X(P_X, \sigma) \;\simeq\; H^0(X, \omega_X),
  \qquad u_\omega \longleftrightarrow \omega,
\]
isomorphisme de groupes, la source étant additive.

La question posée à la fin de la page est celle que les deux suivantes
résolvent : peut-on avoir $u_\omega(c) = c$, si $c$ est cette fois la classe
d'une connexion projective ?\footnote{Le feuillet emploie la lettre $c$ pour
deux objets différents dans le même paragraphe — la classe d'extension dans
$H^1(X,\omega)$ et la connexion projective — et la question finale est
précisément à la charnière. On a distingué les deux ici. Le feuillet mentionne
aussi $P(P^1_{X/S}(-1))$, c'est-à-dire le fibré projectif des parties
principales d'ordre $1$ d'un module inversible : c'est la construction
standard du fibré sous-jacent à une structure projective, et le feuillet ne
la développe pas.}

\pagerange{19}{20}
\section*{Le calcul en coordonnées, et la formule encadrée (pages 19 et 20)}

Soit $\mathcal{U} = \mathfrak{g} = \mathrm{ad}(P)$ le fibré adjoint, de rang
$3$. Le choix d'une uniformisante $t$ le scinde :
\[
  \mathcal{U} \;\simeq\; \omega \oplus \mathcal{O}_X \oplus \omega^{-1},
\]
les trois facteurs étant les espaces de poids $2, 0, -2$ de l'action adjointe
de $\mathfrak{sl}_2$. La coordonnée $t$ définit une connexion projective
$c_{D_t} : \mathcal{U} \to \mathcal{U} \otimes \omega \simeq
\omega^{\otimes 2} \oplus \omega \oplus \mathcal{O}_X$, déterminée par
\[
  c_{D_t}(dt) = -2\,dt, \qquad c_{D_t}(1) = 1, \qquad
  c_{D_t}\big((dt)^{-1}\big) = 0,
\]
les coefficients $-2$ et $0$ étant ceux de l'action adjointe sur les vecteurs
de poids extrême.

Pour $\omega \in \Gamma(X, \omega)$ on pose $u_\omega = \exp \Theta_\omega$,
l'automorphisme unipotent de $\mathcal{U}$ engendré par $\mathrm{ad}(\omega)$,
et l'on fait agir $u_\omega$ sur les connexions par conjugaison,
$u_\omega(c) = u \circ c \circ u^{-1}$. Le calcul de la page 19 — trois
niveaux de crochets, avec $[-\omega, 1] = -2\omega$ — aboutit page 20 au
résultat encadré :
\[
  \boxed{\;u_\omega(c) \;=\; c + d^1_{\omega/S}(\omega)\;}
\]
pour toute connexion projective $c$, où $d^1_{\omega/S}$ est un opérateur
différentiel d'ordre $1$ de $\omega$ vers $\omega^{\otimes 2}$.

C'est la formule structurante du dossier, et il faut voir ce qu'elle dit :
l'ensemble des connexions projectives sur $X$ est un \emph{espace affine} sous
$\Gamma(X, \omega^{\otimes 2})$ — la différence de deux connexions projectives
est une différentielle quadratique — et le groupe $H^0(X, \omega)$ des
automorphismes du couple $(P,\sigma)$ y agit par les translations de l'image
de $d^1_{\omega/S}$. C'est exactement ce que la théorie formelle des
pages 24 à 33 va démontrer abstraitement, sous la forme
$G/G_0 \simeq \Gamma(X, \Omega^{\otimes 2})$.

Les deux corollaires que le feuillet tire sont mutilés au point décisif, et on
n'y touche pas. Ils portent tous deux sur la condition
$d^1_{\omega/S}(\omega) = 0$, avec la mention « a fortiori $d\omega = 0$ » et
la restriction « si $2$ est inversible dans $S$ ». Ce qu'ils visent est clair :
que l'action de $H^0(X,\omega)$ sur les connexions projectives soit libre,
donc que deux connexions projectives distinctes ne donnent pas le même fibré à
structure projective. Ce qui les rend illisibles ne l'est pas moins : la
condition $d\omega = 0$ est vide sur une courbe relative, puisque
$\Omega^2_{X/S} = 0$ et que toute forme globale y est fermée.\footnote{On ne
reconstruit donc ni l'un ni l'autre corollaire. Deux remarques, pour qui
reprendra le feuillet. Premièrement, l'opérateur $d^1_{\omega/S}$ n'est pas la
différentielle extérieure : le feuillet les distingue lui-même, en écrivant
que $d^1_{\omega/S}(\omega) = 0$ entraîne « a fortiori » $d\omega = 0$. C'est
un opérateur d'ordre $1$ attaché à la structure projective, dont $d$ serait
une composante. Deuxièmement, le corollaire a) est illisible juste à
l'endroit de la négation — « $u_\omega(c)$ n'est une connexion projective
[\ldots] ssi » — de sorte qu'on ne sait pas si l'énoncé affirme ou nie. Le
bloc qui suit, où se lisent $M_g(X)$, $M_g^{\mathfrak{z}}(X)$ et
$\mathrm{conn}(X)$ reliés par des flèches, est entièrement biffé de traits
obliques.}

Les dernières lignes de la page 20 ouvrent l'argument que les pages 21 et 22
achèvent : si un tel automorphisme existait, le fibré à structure projective
serait affine, donc isomorphe à un $P(F)$ pour une extension
$0 \to M \to F \to \mathcal{O}_X \to 0$ à structure projective, avec
$E \simeq F \otimes L$. C'est exactement l'hypothèse dont la page 21 part.

\pagerange{21}{22}
\section*{Aucune réduction du groupe structural : la densité de la monodromie (pages 21 et 22)}

\pagerange{21}{21}
\subsection*{Le calcul de degrés}

Reprenons la situation où la page 20 s'est arrêtée. $E$ est de rang $2$ avec
$\det E \simeq \omega_X$, et l'on suppose que le fibré projectif associé se
réduit à un sous-groupe de Borel, c'est-à-dire que $E \simeq F \otimes L$ avec
$F$ muni d'une réduction et $L$ inversible. De
$\det E \simeq \det F \otimes L^{\otimes 2}$, avec $\det F = M$ de degré $0$,
on tire
\[
  2 \deg L = \deg \omega_X = 2g-2, \qquad \text{donc} \qquad \deg L = g-1 .
\]
Considérons alors le composé
\[
  \omega \hookrightarrow E \simeq F \otimes L
  \twoheadrightarrow L .
\]
C'est une section de $L \otimes \omega^{-1}$, module de degré
$(g-1) - (2g-2) = -(g-1) < 0$ dès que $g \geqslant 2$ : elle est donc nulle.
Par conséquent $E \to L$ se factorise par $E/\omega \simeq \mathcal{O}_X$,
d'où $L \simeq \mathcal{O}_X$ et $\deg L = 0$, ce qui contredit
$\deg L = g-1 \neq 0$.

Le groupe structural ne peut donc pas se réduire à un sous-groupe de Borel.
Mieux : il ne peut pas non plus se réduire au normalisateur d'un tore maximal,
car le passage au revêtement double correspondant ramènerait au tore lui-même,
donc à un Borel — et le genre du revêtement est encore $\geqslant 2$, de sorte
que l'argument s'y applique.

\begin{quote}
\textbf{Théorème} (page 21). Soit $X$ une courbe projective lisse de genre
$g \geqslant 2$ sur un corps $k$ de caractéristique nulle. Alors pour toute
connexion projective sur $X$, le groupe structural $H = \mathrm{PGL}_2$ du
fibré principal associé ne se réduit à aucun sous-groupe algébrique $H'$
distinct de $H$.
\end{quote}

\textbf{Une lacune, et comment la combler.} Les deux cas traités
n'épuisent pas les sous-groupes algébriques propres de $\mathrm{PGL}_2$ en
caractéristique $0$. La classification en donne trois familles : ceux qui sont
contenus dans un Borel, ceux qui sont contenus dans le normalisateur d'un tore
maximal, et les trois groupes exceptionnels $A_4$, $S_4$, $A_5$, qui ne sont
dans ni l'un ni l'autre. Ces trois-là se traitent pourtant par le procédé même
que le feuillet emploie pour le normalisateur du tore : une réduction à un
sous-groupe fini $\Gamma$ donne une monodromie d'image finie, donc un
revêtement étale $Y \to X$ de degré $|\Gamma|$ sur lequel la monodromie est
triviale ; la connexion projective se relève à $Y$, le genre de $Y$ est encore
$\geqslant 2$ par Riemann-Hurwitz, et le groupe structural y est trivial, donc
contenu dans un Borel — ce que le premier cas exclut.\footnote{Le feuillet
écrit « Le [\ldots] argument montre : » et l'adjectif est illisible, juste à
l'endroit où il aurait dit s'il considérait le cas général comme acquis. La
page 21 est, avec les pages 6, 7 et 11, l'une des quatre plus abîmées du
dossier : le nom du groupe y est écrit d'un seul jet et les lettres n'en sont
pas assurées, et un symbole isolé qualifiant les fibrés n'a pas été lu du
tout. Le théorème est vrai tel qu'il est énoncé ; ce que cette note ajoute est
le troisième cas, que le feuillet ne traite pas.}

\pagerange{22}{22}
\subsection*{La densité de Zariski, et où la situer}

Les trois lignes de la page 22 tirent la conséquence analytique. Sur
$k = \mathbf{C}$, l'homomorphisme de monodromie
\[
  \pi_1(X^{\mathrm{an}}, x) \longrightarrow H(\mathbf{C})
\]
a une image Zariski-dense dans $H_{\mathbf{C}}$ : car une image non dense
aurait pour adhérence un sous-groupe algébrique propre, auquel le fibré se
réduirait.

Il vaut la peine de situer cet énoncé, puisque c'est le genre de résultat dont
un lecteur moderne voudra savoir où il se place. Ce que le dossier démontre
est, dans le vocabulaire d'aujourd'hui, l'\emph{irréductibilité de la
monodromie d'une structure projective} sur une courbe de genre $\geqslant 2$ :
un $\mathfrak{sl}_2$-oper, au sens de Beilinson et Drinfeld, a une monodromie
Zariski-dense. C'est un énoncé standard, et la démonstration standard est
exactement celle du feuillet — un sous-fibré en droites invariant aurait le
degré $g-1$, incompatible avec la position du fibré sous-jacent dans sa
filtration de Harder-Narasimhan. Du côté analytique, la même chose se dit :
une structure projective sur une surface de Riemann compacte de genre au moins
$2$ a une monodromie non élémentaire, ce qu'on trouve chez Gunning dans les
années soixante.

Cette édition n'affirme aucune priorité, et n'a pas fait de recherche
bibliographique exhaustive. Ce qu'elle constate est que le résultat est
aujourd'hui connu, qu'il est ici obtenu par l'argument qui est resté le bon,
et qu'il est, dans le dossier lui-même, la pièce dont les pages 14 à 16 ont
besoin sans le dire.

\pagerange{24}{33}
\section*{La théorie formelle : rigidité d'une algèbre à connexion (pages 24 à 33)}

Le dernier tiers du dossier change de sujet. Il n'y est plus question de
courbes, ni de genre, ni de monodromie : seulement d'une algèbre complète et
d'une dérivation. C'est la démonstration technique de ce que les pages 7 à 9
avaient posé sans preuve, et de ce que la page 20 avait obtenu par le calcul.
Il est repris cinq ou six fois ; on suit ici l'ordre logique, en signalant à
chaque fois de quelle mouture on parle.

\pagerange{32}{33}
\subsection*{Le cadre, tel que les pages 32 et 33 le posent}

C'est la dernière mouture qui est la plus propre ; on commence par elle.

Soient $X/S$ un schéma relatif, et $P$ une $\mathcal{O}_X$-Algèbre augmentée,
séparée et complète pour la topologie $J$-adique, où $J$ est l'idéal
d'augmentation, avec $J/J^2 = \Omega$ inversible. Soit
\[
  C : P \longrightarrow P \,\widehat{\otimes}\, \Omega^1_{X/S}
\]
une connexion relative à $S$, compatible à la structure d'algèbre :
\[
  \begin{cases}
    C(fg) = f\,C(g) + g\,C(f), & f,g \text{ sections de } P, \\
    C(\lambda \cdot 1) = 1 \otimes d\lambda, &
    \lambda \text{ section de } \mathcal{O}_X .
  \end{cases}
\]

Localement sur $X$, on choisit un isomorphisme d'algèbres augmentées
$P \simeq \widehat{\operatorname{Sym}}(\Omega)$ et l'on écrit $C$ dans ce
modèle. La restriction de $C$ à $\Omega = \operatorname{Sym}^1$ donne
\[
  C|_\Omega : \Omega \longrightarrow
  \widehat{\operatorname{Sym}}(\Omega) \,\widehat{\otimes}\, \Omega^1_{X/S}
  = \prod_{i \geqslant 0} \operatorname{Sym}^i(\Omega)
  \otimes \Omega^1_{X/S},
\]
dont on note $(\varpi_i)_{i \geqslant 0}$ les composantes. Leur nature se lit
sur la règle de Leibniz : pour $i \neq 1$ la composante $\varpi_i$ est
$\mathcal{O}$-linéaire, donc une section globale,
\[
  \varpi_i \in \Gamma\big( \Omega^{\otimes(i-1)} \otimes \Omega^1_{X/S}
  \big) \quad (i \neq 1),
\]
tandis que $\varpi_1$ est une \emph{connexion sur $\Omega$} et non une
section.\footnote{La distinction est celle que le feuillet écrit en gras :
« les $\varpi_i$ pour $i \neq 1$ sont linéaires ». Elle est forcée par
$C(\lambda\omega) = \lambda\,C(\omega) + \omega\,d\lambda$ de la page 27 : le
terme $\omega\,d\lambda$ n'affecte que la composante de degré $1$. Le symbole
$\varpi$ est écrit d'un seul trait bouclé sur les feuillets et le
transcripteur le fixe ainsi dans toute la suite ; on garde sa convention.}

Le symbole de $C$ est la donnée essentielle. On a
$C(J^i) \subset J^{i-1} \,\widehat{\otimes}\, \Omega^1_{X/S}$, d'où pour tout
$i \geqslant 1$ un homomorphisme $\mathcal{O}$-linéaire sur les gradués
\[
  \operatorname{gr}^i = J^i/J^{i+1} \simeq \Omega^{\otimes i}
  \xrightarrow{\ \widetilde{c}\ }
  \Omega^{\otimes(i-1)} \otimes \Omega^1_{X/S},
\]
et cet homomorphisme est la multiplication par un seul et même
\[
  \varpi_0 \in \Gamma\big( \Omega^{-1} \otimes \Omega^1_{X/S} \big)
  = \mathrm{Hom}\big( \Omega, \Omega^1_{X/S} \big),
\]
\emph{indépendant de $i$}. On dit que $C$ est \textbf{régulière} si $\varpi_0$
l'est, c'est-à-dire si
\[
  \varpi_0 : \Omega \xrightarrow{\ \sim\ } \Omega^1_{X/S} .
\]
C'est la condition qui fait de $P$ une algèbre de jets et non une algèbre
quelconque : c'est l'identification $(*)$ de la page 7, retrouvée comme une
propriété de la connexion plutôt que comme une donnée.

\pagerange{25}{26}
\subsection*{Le théorème de rigidité (pages 25 et 26)}

Voici l'énoncé, dans la mouture des pages 25 et 26, qui est la plus
articulée. On se place en caractéristique $0$.

\begin{quote}
\textbf{Théorème} (page 25). Soit $P$ comme ci-dessus, munie d'une connexion
$C$ régulière. Alors il existe un couple $(P,\sigma)$ — c'est-à-dire une
extension $\mathcal{E} : 0 \to \Omega \to E \to \mathcal{O}_S \to 0$, et donc
un modèle $\widehat{S} = \widehat{\operatorname{Sym}}(\Omega)$ muni de sa
connexion canonique $\widehat{c}$ — et un isomorphisme d'algèbres augmentées
$u : \widehat{S} \xrightarrow{\ \sim\ } P$ tel que $u(C) = \widehat{c}$.
\end{quote}

Le théorème se démontre en deux temps, et c'est cette décomposition qui porte
tout le contenu :

\begin{itemize}
\item[a)] La donnée d'un isomorphisme $u$ \emph{à l'ordre $2$} détermine déjà
$(P,\sigma)$, c'est-à-dire l'extension $\mathcal{E}$, à isomorphisme unique
près, par $\Omega = J/J^2$ et $E$ construit sur $J/J^3$.
\item[b)] $(P,\sigma)$ étant ainsi déterminé, pour tout isomorphisme $u_3$
\emph{à l'ordre $3$}, $\widehat{S}/\widehat{J}^{\,4} \to P/J^4$, prolongeant
celui d'ordre $2$, il existe un isomorphisme $u$ et un seul qui le prolonge et
tel que $u(C)$ soit de la forme $\widehat{c}$.
\end{itemize}

Autrement dit : rien n'est libre au-delà de l'ordre $3$. Toute l'ambiguïté est
concentrée au passage de l'ordre $2$ à l'ordre $3$, et la page 26 l'identifie
en donnant une correspondance biunivoque entre trois
ensembles.\footnote{L'item b) du théorème s'arrête au bas du feuillet 25, la
phrase inachevée, et se termine en tête de la page 26 : « il existe un iso. et
un seul $u$ tel [qu'il prolonge] $u_2$, tel que $u(\widehat{c}) = C$ ». Dans
l'item a) du feuillet, le signe entre $\Omega$ et $(J/J^3)$ est lu $\otimes$
sans certitude par le transcripteur ; on a préféré écrire « $E$ construit sur
$J/J^3$ », qui est ce que la construction demande, plutôt que de reproduire un
produit tensoriel dont la lecture n'est pas assurée.}

\begin{itemize}
\item[a)] les solutions du problème, à isomorphisme près — la catégorie des
solutions étant \emph{rigide}, c'est-à-dire sans automorphismes non triviaux ;
\item[b)] $(P,\sigma)$ et $u_2$ étant fixés, l'ensemble des
pseudo-$u_3$ qui prolongent $u_2$, lequel est un \emph{torseur sous}
$\Gamma(X, \Omega^{\otimes 2})$ ;
\item[c)] l'ensemble $\mathrm{Conn}^*(P)$ des connexions régulières sur $P$
dont l'image est de la forme $\widehat{c}$.
\end{itemize}

C'est le résultat que tout le dossier attendait. La fibre de
$M_g^{\mathfrak{z}} \to M_g$ est un torseur sous
$\Gamma(X, \Omega^{\otimes 2}) = H^0(X, \omega^{\otimes 2})$, donc de
dimension $3g-3$ ; c'est le contenu abstrait de la formule encadrée de la
page 20.\footnote{L'item c) est presque entièrement perdu sur le feuillet :
on y lit « l'ens. des connexions $C$ sur $P$ telles que $\widehat{c}$ soit,
dans [\ldots], [donnée] de $C$ [\ldots] sur $\mathrm{Conn}^*(P)$ », et ni le
complément de « dans », ni le verbe, ne se laissent lire. Ce qui est sûr est
que le troisième terme de la correspondance est $\mathrm{Conn}^*(P)$, que les
pages 28 et 33 nomment et définissent, et que le petit diagramme du bas de la
page 26 confirme : $\mathrm{Sol} \xrightarrow{\sim} \mathrm{Iso}_3$ avec une
oblique vers $\mathrm{Conn}^*P$. On n'en a pas reconstruit davantage. Le bas
de la page 26 est du brouillon de calcul : une quinzaine d'égalités jetées
côte à côte, la plupart biffées, où se lit l'essentiel de l'argument de la
page 29 sous forme d'aide-mémoire.}

\pagerange{27}{28}
\subsection*{Le calcul du changement de paramètre (pages 27, 28 et 30)}

Les trois feuillets calculent la même chose, la page 30 étant la plus
complète. Le point est de vérifier, à la main, que le torseur annoncé est bien
celui-là et que rien ne bouge au-delà.

On écrit, pour un générateur local $\omega_0$ de $\Omega$,
\[
  C(\omega_0) = \sum_{i \geqslant 0} \omega_0^{\,i} \otimes \varpi_i,
\]
et l'on change de générateur : $\omega_0 \mapsto \omega_0\,u$ avec
\[
  u = 1 + \sum_{i \geqslant 2} \alpha_i\, \omega_0^{\,i} .
\]
Le feuillet note $\omega_0^{\,u}$ le produit $\omega_0 u$ ; le $u$ surélevé
est un \emph{facteur, non un exposant}.\footnote{C'est la page 30 qui le dit
en toutes lettres, en écrivant
$C(\omega_0^{\,u}) = C\big(\omega_0(1 + \sum_{i\geqslant 2}
\alpha_i\omega_0^{\,i})\big)$, ce qui identifie $\omega_0^{\,u}$ au produit
$\omega_0 u$ pour ce même $u$ qui figure aux dénominateurs des formules de la
page 28. La transcription conserve la notation telle quelle et le signale ;
on la conserve ici aussi, avec l'avertissement. Sans lui, les formules de la
page 28 sont illisibles.}

La règle de Leibniz donne alors
\[
  C(\omega_0^{\,u}) = C(\omega_0)\Big( 1 + \sum_{i \geqslant 2}
  (i+1)\,\alpha_i\, \omega_0^{\,i} \Big) + \text{(termes en } d\alpha_i),
\]
le coefficient $(i+1)$ venant de la somme du facteur $\alpha_i\omega_0^i$ et
de la dérivation $i\,\alpha_i\omega_0^{i-1}C(\omega_0)$ du produit. Ramenant
par $u^{-1}$ et travaillant modulo $\Omega^{n+1}$ avec $\alpha_i = 0$ pour
$i < n$, il reste
\[
  C(\omega_0^{\,u})^{u^{-1}} = C(\omega_0)
  + (n+1)\,\alpha_n\, \omega_0^{\,n}\, \varpi_0,
\]
de sorte que la condition pour que le changement de paramètre préserve la
forme de la connexion est
\[
  (n+1)\,\alpha_n = 0 ,
\]
c'est-à-dire, en caractéristique $0$, $\alpha_n = 0$. C'est la rigidité, sous
sa forme la plus élémentaire : au-delà du terme de degré $2$, aucun changement
de paramètre n'est disponible. Le feuillet vérifie les cas $n = 2$ et
$n = 3$ séparément avant d'écrire le cas général.\footnote{Le premier terme du
développement de $C(\omega)$ à la page 30 est transcrit
$C(\omega) = \omega\,\varpi_{-1} + \omega\,\varpi_0 + \ldots$, avec la note
que ce premier terme n'est pas lu avec certitude. Il ne peut être ce qui est
écrit : la page 27 pose $C(\omega_0) = \sum_{i \geqslant 0} \omega_0^i \otimes
\varpi_i$, où l'indice part de $0$, et un $\varpi_{-1}$ n'y a pas de place.
Cette édition n'en propose pas de lecture et n'utilise que la formule de la
page 27. La moitié inférieure des pages 27 et 30 est du brouillon, en crayon
très pâle, largement biffé ; on n'en a rien tiré.}

\pagerange{29}{31}
\subsection*{La structure du groupe $G$ (pages 29 et 31)}

Soit $G$ le groupe des automorphismes de $\widehat{S}$ comme algèbre
augmentée qui sont l'identité modulo $J^3$, et $G_0 \subset G$ le sous-groupe
de ceux qui sont l'identité modulo $J^4$. Un élément de $G$ est déterminé par
l'image de $\Omega$, soit $\omega \mapsto \omega + q_3(\omega) + \ldots$ avec
$q_i : \Omega \to \Omega^{\otimes i}$ ; être l'identité modulo $J^3$ annule
$q_2$, l'être modulo $J^4$ annule en outre $q_3$. D'où
\[
  G/G_0 \;\simeq\; \mathrm{Hom}\big(\Omega, \Omega^{\otimes 3}\big)
  \;\simeq\; \Gamma\big(\Omega^{\otimes 2}\big),
\]
le dernier isomorphisme parce que $\Omega$ est inversible.\footnote{Le calcul
de $G/G_0$ est esquissé par le feuillet, qui écrit directement
« donc $G/G_0 \simeq \Gamma(\Omega^{\otimes 2})$ » sans le justifier ;
l'argument par les $q_i$ est de cette édition, et c'est lui qui explique que
ce soit $\Omega^{\otimes 2}$ et non $\Omega^{\otimes 3}$ qui apparaisse. C'est
aussi ce qui rend la fibre de $M_g^{\mathfrak{z}} \to M_g$ de dimension
$3g-3$ et non $4g-4$.}

On pose, pour une connexion $c$ sur $P$,
\[
  G(c) = \{\, h \in G \ \mid\ h \cdot \widehat{c}
  \text{ est de la forme } \widehat{c}\,' \,\},
\]
qui est un sous-groupe de $G$. Le théorème de structure est alors :
\[
  G_0 \times G(c) \longrightarrow G, \qquad (g,h) \longmapsto g h
\]
est bijectif, d'où $G(c) \xrightarrow{\ \sim\ } G/G_0 \simeq
\Gamma(X, \Omega^{\otimes 2})$.

La démonstration tient en deux lignes de chaque côté. Injectivité : si
$gh = g'h'$, alors $(g'^{-1}g)h = h'$ ; or $h\widehat{c}$ et $h'\widehat{c}$
sont tous deux de la forme $\widehat{c}\,'$, et le corollaire du haut de la
page 29 — un élément de $G_0$ qui envoie $\widehat{c}$ sur un $\widehat{c}\,'$
est l'identité — donne $g'^{-1}g = 1$, d'où $g = g'$ puis $h = h'$.
Surjectivité : pour $g \in G$, la proposition de la page 28 écrit
$g\widehat{c} = k\,\widehat{c}_0$ avec $k \in G_0$, d'où $k^{-1}g \in G(c)$ et
$g = k(k^{-1}g)$.

La page 31 reprend exactement cette démonstration en la dépouillant de toute
géométrie : un groupe $G$, un sous-groupe $G_0$, un $G$-ensemble $C$, une
partie $C_0$ telle que $G_0 \times C_0 \xrightarrow{\ \sim\ } C$, et
$G(c) = \{g \in G \mid g\cdot c \in C_0\}$ pour $c \in C_0$ ; la conclusion est
la même, et l'énoncé n'a plus rien d'algébrique. C'est le même argument vu
deux fois, et il vaut la peine de garder les deux : la seconde version montre
que le résultat ne doit rien aux algèbres, seulement à l'action
simplement transitive.\footnote{La démonstration de la page 31 emploie la
lettre $h$ pour trois objets différents dans la même phrase — « soit
$h \in G$, alors $hc$ est de la forme $h \cdot c'$ [\ldots] donc $h^{-1}k$
[\ldots] donc $h = hg$ » — et la transcription reproduit fidèlement la
collision. On a récrit la démonstration avec des lettres distinctes ; c'est la
version de la page 29, avec $k$, qui est correcte telle quelle, et c'est elle
qu'on a suivie.}

Le corollaire que la page 31 tire du lemme — l'intersection d'une orbite de
$G$ dans $C$ avec $C_0$ s'identifie à $G_0 \backslash G$ — est la même chose
dite du côté des orbites. Son énoncé est mutilé sur le feuillet, et on ne le
mentionne que pour mémoire.

\pagerange{33}{33}
\subsection*{L'énoncé final}

La page 33 referme, avec ses hypothèses au complet. En caractéristique
nulle, pour $P = \widehat{S} = \mathcal{O}_{\widehat{P}}$ :
\[
  \mathrm{Conn}^*(P) \times G_0 \longrightarrow
  \mathrm{Conn}^*(\widehat{S}), \qquad (c,g) \longmapsto g \cdot \widehat{c}
\]
est bijective, où $\mathrm{Conn}^*$ désigne les connexions \emph{régulières} ;
et pour $c \in \mathrm{Conn}^*(P)$,
\[
  G(c) \;\xrightarrow{\ \sim\ }\; G/G_0 \;\simeq\;
  \Gamma\big(X, \Omega^{\otimes 2}\big).
\]

C'est le dernier mot du dossier, et c'est aussi le plus important : l'ensemble
des structures projectives sur une courbe fixée est un torseur sous les
différentielles quadratiques.

\pagerange{1}{33}
\section*{Ce que le dossier cherche, et ce qui lui manque}

Reste à dire ce que ces trente-trois pages veulent, puisque le titre de la
page 1 est de lui et annonce l'uniformisation.

\pagerange{1}{33}
\subsection*{Ce qui s'articule}

Trois choses se rejoignent, et deux d'entre elles ne sont jamais rapprochées
par le dossier.

\begin{itemize}
\item \textbf{La théorie formelle des pages 24 à 33 est la démonstration de ce
que les pages 7 à 9 posaient.} La définition d'une connexion projective exige
un isomorphisme $P^n_{X/S} \simeq \mathcal{O}_{\widehat{P}}/\widehat{J}^{n+1}$
pour tout $n$ ; le théorème de la page 25 dit que cet isomorphisme existe, et
qu'il est unique une fois fixé à l'ordre $3$. Sans lui, la définition des
pages 7 à 9 demanderait une infinité de données et ne serait pas maniable.
Cette articulation-là, le dossier la fait : le vocabulaire est le même, les
notations $\widehat{S}$, $\widehat{J}$, $\widehat{c}$ sont celles de la
page 3.
\item \textbf{La théorie formelle calcule la fibre que les pages 14 à 18
comptent.} L'ambiguïté à l'ordre $3$ est un torseur sous
$\Gamma(X,\Omega^{\otimes 2})$, de dimension $3g-3$ : c'est exactement la
fibre de $M_g^{\mathfrak{z}} \to M_g$, et avec $\dim M_g = 3g-3$ cela donne
$\dim M_g^{\mathfrak{z}} = 6g-6$, qui est le nombre que la page 15 obtient par
un tout autre chemin, la caractéristique d'Euler-Poincaré de de Rham du fibré
adjoint. Deux calculs indépendants, le même résultat. Le dossier ne les
confronte jamais — ils sont séparés par onze feuillets et un changement de
sujet — et c'est pourtant leur accord qui rend l'anomalie de la page 15
lisible comme une inadvertance et non comme un désaccord.
\item \textbf{L'argument de densité des pages 21 et 22 fournit l'hypothèse
dont la page 14 a besoin.} Le critère de non-dégénérescence de la page 14
n'est complet que si les automorphismes finis sont exclus aussi bien que les
infinitésimaux ; la densité Zariski de la monodromie les exclut tous. Cette
articulation-là non plus n'est pas faite par le dossier, et elle est de cette
édition.
\end{itemize}

\pagerange{1}{33}
\subsection*{Ce qui ne s'articule pas}

Il faut être franc sur le reste, parce que le titre promet davantage que ce
que les pages livrent.

\begin{itemize}
\item \textbf{Il n'y a aucune connexion projective canonique dans le dossier.}
Le mot est au titre ; la page 8 pose la question — une connexion projective
est-elle canoniquement déterminée par $X/S$, sans donnée supplémentaire ? — et
la laisse, avec un point d'interrogation. Tout ce que le dossier établit va
dans l'autre sens : les connexions projectives forment un espace affine de
dimension $3g-3$, sans point privilégié. Un point privilégié ne peut venir que
d'ailleurs — de l'uniformisation, précisément, qui en fournit un sur
$\mathbf{C}$ mais pas algébriquement.
\item \textbf{L'uniformisation elle-même n'est pas atteinte.} Ce qu'on en
trouve est le plongement de l'espace de Teichmüller dans la variété des
caractères, page 17, et la caractérisation de son image par l'existence d'un
domaine invariant à quotient compact. La caractérisation est juste et elle est
la bonne ; mais le feuillet l'introduit par « on suppose », et rien n'est
démontré. La question qu'il pose ensuite — l'injectivité de l'application de
monodromie — est restée une vraie question.
\item \textbf{Les deux massifs ne se parlent pas.} Les pages 1 à 22 sont de la
géométrie des courbes ; les pages 24 à 33 sont de l'algèbre formelle
au-dessus d'un $X/S$ quelconque, où le genre n'apparaît jamais et où le mot
« courbe » n'est pas prononcé. Rien dans le second massif ne renvoie au
premier, et l'unique pont — que $\Gamma(X,\Omega^{\otimes 2})$ du dernier
feuillet est le $H^0(X, \omega^{\otimes 2})$ des pages 12 et 17 — n'est
jamais franchi sur le papier.
\end{itemize}

Ce que le dossier est, finalement : l'appareil algébrique complet d'une
théorie des structures projectives sur les courbes, avec un théorème net au
milieu, et sans le théorème que le titre annonce. C'est un état de chantier, et
c'est ce qui le rend lisible — on y voit ce qu'il fallait construire avant de
pouvoir poser la question.

\end{document}
